\documentclass[a4paper]{article}
% Español (México) con XeLaTeX: fontspec para fuentes Unicode, polyglossia
% para la división silábica y los nombres del documento en la variante mexicana.
\usepackage{fontspec}
\usepackage{polyglossia}
\setdefaultlanguage[variant=mexican]{spanish}
% Los nombres chinos van como «pinyin (original)»: xeCJK da a los caracteres
% Han su propia fuente; sin ella XeLaTeX los omite con un aviso.
% FandolSong no cubre algunos caracteres de nombres (U+73FA); WenQuanYi Zen Hei sí.
\usepackage[AutoFallBack=true]{xeCJK}
\setCJKmainfont{FandolSong-Regular.otf}
\setCJKfallbackfamilyfont{\CJKrmdefault}{WenQuanYi Zen Hei}
% polyglossia no traduce los nombres de listings.
\AtBeginDocument{\providecommand{\lstlistingname}{}\renewcommand{\lstlistingname}{Listado}%
  \providecommand{\lstlistlistingname}{}\renewcommand{\lstlistlistingname}{Índice de listados}}
\usepackage{amsmath,amssymb}
\usepackage{graphicx}
\usepackage[margin=2.35cm]{geometry}
\usepackage[most]{tcolorbox}
\usepackage{etoolbox}
\usepackage{listings}
\usepackage{xcolor}
\usepackage{hyperref}
\usepackage{booktabs,longtable,tabularx,array}
\usepackage{subcaption}
\usepackage{float}
\usepackage{enumitem}
\usepackage{microtype}
\usepackage{fancyhdr}
\usepackage{needspace}

\definecolor{kbBack}{HTML}{F2F6FA}
\definecolor{kbFrame}{HTML}{9DB4CC}
\definecolor{kbTitle}{HTML}{52708F}
\definecolor{ibBack}{HTML}{FBF5E7}
\definecolor{ibFrame}{HTML}{C9A961}
\definecolor{ibTitle}{HTML}{7D5F1E}
\definecolor{wbBack}{HTML}{FCF2F0}
\definecolor{wbFrame}{HTML}{D6A096}
\definecolor{wbTitle}{HTML}{9E4F43}
\definecolor{dbBack}{HTML}{F4F7F3}
\definecolor{dbFrame}{HTML}{AEC2AC}
\definecolor{dbTitle}{HTML}{5F7F64}

\tcbset{softbox/.style={
  enhanced,breakable,boxrule=0.8pt,arc=2pt,left=3mm,right=3mm,
  top=2mm,bottom=2mm,coltitle=white,fonttitle=\bfseries,
  toptitle=0.6mm,bottomtitle=0.6mm
}}
\newtcolorbox{knowledgebox}[1]{softbox,colback=kbBack,colframe=kbFrame,colbacktitle=kbTitle,title={#1}}
\newtcolorbox{importantbox}[1]{softbox,colback=ibBack,colframe=ibFrame,colbacktitle=ibTitle,title={#1}}
\newtcolorbox{warningbox}[1]{softbox,colback=wbBack,colframe=wbFrame,colbacktitle=wbTitle,title={#1}}
\newtcolorbox{dialoguebox}[1]{softbox,colback=dbBack,colframe=dbFrame,colbacktitle=dbTitle,title={#1}}

\lstset{
  language=Python,basicstyle=\ttfamily\small,
  keywordstyle=\color{kbTitle}\bfseries,stringstyle=\color{wbTitle},
  commentstyle=\color{dbTitle},breaklines=true,frame=single,numbers=left,
  numberstyle=\tiny\color{gray},captionpos=b,columns=fullflexible,
  keepspaces=true,showstringspaces=false,extendedchars=true
}
\BeforeBeginEnvironment{lstlisting}{\Needspace{12\baselineskip}}

\hypersetup{colorlinks=true,linkcolor=kbTitle,urlcolor=kbTitle,citecolor=wbTitle}
\setlist{itemsep=0.25em,topsep=0.4em}
\setlength{\parindent}{2em}
\setlength{\parskip}{0.25em}
\newcolumntype{L}[1]{>{\raggedright\arraybackslash}p{#1}}

\providecommand{\notetitle}{Notas del curso: [tema del curso]}
\providecommand{\noteauthors}{Elaboradas a partir de los materiales públicos de Stanford CS25}
\providecommand{\notedate}{Versión reescrita en 2026}
\providecommand{\videochannel}{Stanford Online}
\providecommand{\videopublishdate}{2022}
\providecommand{\videoduration}{[duración del video]}
\providecommand{\videourl}{}
\providecommand{\videocoverpath}{cover.jpg}
\providecommand{\coursesite}{https://web.stanford.edu/class/cs25/}
\providecommand{\repourl}{https://github.com/hqhq1025/ai-course-notes}
\providecommand{\skillversion}{wdkns/wdkns-skills@39f1a04}

\newcommand{\figsource}[1]{\par\vspace{-0.3em}{\footnotesize\color{gray}\emph{Fuente: #1}}\par}
\newcommand{\teachervoice}[1]{\begin{knowledgebox}{Nota de clase}#1\end{knowledgebox}}
\newcommand{\term}[1]{\textbf{#1}}

\pagestyle{fancy}
\fancyhf{}
\fancyfoot[C]{\thepage}
\fancyfoot[R]{\footnotesize\color{gray}\href{\repourl}{AI Course Notes}}
\renewcommand{\headrulewidth}{0pt}
\renewcommand{\footrulewidth}{0pt}

\newcommand{\makecscover}{
\begin{titlepage}
\centering
\vspace*{0.7cm}
{\Large Stanford CS25 · Transformers United\par}
\vspace{0.8cm}
{\huge\bfseries \notetitle\par}
\vspace{0.6cm}
{\large \noteauthors\par}
\vspace{0.25cm}
{\large \notedate\par}
\vspace{0.9cm}
\ifdefempty{\videocoverpath}{}{
  \includegraphics[width=0.86\textwidth,height=0.43\textheight,keepaspectratio]{\videocoverpath}\par
}
\vfill
\begin{tcolorbox}[width=0.92\textwidth,colback=black!2!white,colframe=black!45,boxrule=0.7pt,arc=2pt]
\textbf{Autor o canal del video}: \videochannel\par
\textbf{Fecha de publicación}: \videopublishdate\par
\textbf{Duración del video}: \videoduration\par
\textbf{Sitio del curso}: \href{\coursesite}{\nolinkurl{\coursesite}}\par
\ifdefempty{\videourl}{}{\textbf{Enlace del video}: \href{\videourl}{\nolinkurl{\videourl}}\par}
\textbf{Normas de elaboración}: \skillversion\ + las normas source-first / teacher-voice / visual-QA de este repositorio
\end{tcolorbox}
\end{titlepage}
\tableofcontents
\newpage
}

\usepackage{multicol}

\renewcommand{\notetitle}{Lección 01: Panorama del Transformer — desde los paradigmas de aprendizaje hasta las arquitecturas de próxima generación}
\renewcommand{\noteauthors}{Basado en las conferencias oficiales de Steven Feng y Karan Singh, el deck de 156 diapositivas y subtítulos manuales}
\renewcommand{\notedate}{Apuntes CS25 V6 · 2026 source-first}
\renewcommand{\videochannel}{Stanford Online}
\renewcommand{\videopublishdate}{Clase 2026-04-02; subida 2026-04-22}
\renewcommand{\videoduration}{1:16:46}
\renewcommand{\videourl}{https://www.youtube.com/watch?v=bHSDPgZYie0}
\renewcommand{\videocoverpath}{cover.jpg}

\newcommand{\lecturefigure}[3]{%
\begin{figure}[H]
  \centering
  \includegraphics[width=0.97\textwidth,height=0.49\textheight,keepaspectratio]{slides-images/#1}
  \caption{#2}
  \figsource{#3}
\end{figure}
}

\let\standardtableofcontents\tableofcontents
\renewcommand{\tableofcontents}{%
  \begingroup
  \small
  \setlength{\parskip}{0pt}
  \standardtableofcontents
  \endgroup
}

\begin{document}

\makecscover

\section{De características manuales a auto-supervisión: ¿qué es lo que realmente aprenden los modelos?}

Esta clase de apertura del curso V6 no simplifica el Transformer reduciéndolo únicamente a una fórmula de atención. Los dos ponentes primero rastrean cómo la señal de aprendizaje evolucionó desde características manuales y etiquetas supervisadas hacia objetivos de auto-supervisión, y luego conectan datos, post-entrenamiento, agentes, multimodalidad, memoria y alineación en un mismo mapa del sistema. Al leer estas notas, debe mantenerse preguntando constantemente tres cuestiones: ¿de dónde proviene la información?, ¿qué recompensan los gradientes? y, ¿qué mecanismos externos al modelo añaden los sistemas una vez desplegados?

\subsection{Promesa del curso: mecanismo, aplicación y límites deben examinarse conjuntamente}

Los puntos clave iniciales sitúan la atención, la escalabilidad, las aplicaciones interdominio, los ponentes de vanguardia y los problemas abiertos en la misma página. Este orden define el principio organizador de las notas: primero explicar el mecanismo, luego contrastar con la evidencia y finalmente explicitar los límites. Conocer solo la forma tensorial de un bloque es insuficiente para explicar por qué un modelo tiene éxito sobre un determinado conjunto de datos; ver únicamente una demostración exitosa tampoco basta para juzgar si es fiable en nuevas distribuciones, contextos largos o flujos de trabajo reales.

\lecturefigure{slide-011.jpg}{Objetivos del curso: desde el mecanismo del Transformer hasta aplicaciones de vanguardia, limitaciones y problemas abiertos}{CS25 V6 oficial deck, p.~11}

\teachervoice{00:06:12--00:07:10，el docente enfatiza que el curso no se limita a explicar la atención, sino que busca que los estudiantes accedan a aplicaciones de vanguardia, limitaciones clave y problemas abiertos en los que aún pueden contribuir. Nota de clase: cada método posterior debe responder a "¿qué capa del sistema ha cambiado?".}

\begin{importantbox}{Las seis capas de interfaz de los sistemas Transformer modernos}
\begin{enumerate}
  \item \textbf{Representación}: cómo la entrada original se convierte en tokens y vectores;
  \item \textbf{Arquitectura}: cómo fluye la información entre secuencias, capas y modalidades;
  \item \textbf{Pre-entrenamiento}: qué capacidades por defecto moldean la distribución de datos y la función objetivo;
  \item \textbf{Post-entrenamiento}: cómo la supervisión, las preferencias, la búsqueda y el verificador modifican el comportamiento;
  \item \textbf{Bucle del sistema}: cómo la recuperación, las herramientas, la memoria y la retroalimentación ambiental ingresan al ciclo cerrado;
  \item \textbf{Gobernanza/evidencia}: cómo se definen la evaluación, la explicabilidad, la alineación y los límites de despliegue.
\end{enumerate}
Un nuevo término que no pueda aterrizar en alguna de estas interfaces generalmente no ha sido verdaderamente explicado.
\end{importantbox}

\subsection{Tres transformaciones: ingeniería de características, representación supervisada y representación auto-supervisada}

El aprendizaje automático temprano dependía de que los investigadores diseñaran manualmente las características, para luego que un modelo relativamente superficial las mapeara a etiquetas. El aprendizaje profundo delegó el aprendizaje de la representación a la red, pero aún exigía abundantes etiquetas manuales. El aprendizaje auto-supervisado continuó desplazando la fuente de supervisión: las etiquetas ya no provienen de categorías humanas, sino que se construyen a partir de la entrada misma, como ocultar tokens, predecir el siguiente paso, restaurar imágenes perturbadas o distinguir muestras coincidentes de no coincidentes. La diferencia clave entre los tres paradigmas no es "qué tan profunda es la red", sino \emph{quién define la tarea de aprendizaje y cuánta estructura original puede ver el modelo}.

\lecturefigure{slide-013.jpg}{Aprendizaje automático temprano: características manuales, optimizador y costosas etiquetas manuales}{CS25 V6 oficial deck, p.~13}

\lecturefigure{slide-014.jpg}{Profundización supervisada: la red comienza a asumir el aprendizaje de representación, pero sigue dependiendo de etiquetas}{CS25 V6 oficial deck, p.~14}

De las diapositivas 13 a la 16 se debe leer según el cuello de botella de información: las características manuales comprimen los datos originales en dimensiones que el investigador considera útiles; la red supervisada amplía el alcance de la representación aprendible; y el entrenamiento end-to-end reduce aún más las interfaces intermedias humanas. El costo también cambia: aumenta la libertad del modelo, lo que requiere más datos y cómputo; si la distribución de entrenamiento presenta sesgos, la red aprenderá esos sesgos junto con los patrones válidos.

\lecturefigure{slide-015.jpg}{Estado transicional del aprendizaje supervisado: los datos originales aún pasan primero por una representación manual}{CS25 V6 oficial deck, p.~15}

\lecturefigure{slide-016.jpg}{Aprendizaje supervisado end-to-end: la entrada original ingresa directamente a la red profunda}{CS25 V6 oficial deck, p.~16}

\begin{knowledgebox}{Lectura de la figura: por qué valen las cuatro páginas progresivas}
Primero compare quién construye la representación, luego desde dónde proviene la señal de supervisión. En la diapositiva 13, los humanos definen las características; en las diapositivas 14--15 coexisten la red y la interfaz manual; en la diapositiva 16, la red consume directamente los datos brutos. La figura respalda el desplazamiento interno gradual de la "interfaz de aprendizaje de representación", pero no prueba que el conocimiento humano carezca de valor. Los sistemas modernos aún inyectan abundante diseño humano mediante tokenizadores, filtrado de datos, priorización arquitectónica, pérdida y APIs de herramientas.
\end{knowledgebox}

El objetivo del aprendizaje auto-supervisado puede abstraerse como construir una parte predecible $y=g(x)$ y un contexto visible $\tilde{x}=c(x)$ a partir de una muestra $x$, para luego optimizar
\[
\mathcal{L}_{\mathrm{SSL}}(\theta)
=-\log p_{\theta}\!\left(y\mid \tilde{x}\right).
\]
$g$ representa la regla para construir el objetivo, como el siguiente token o un parche oculto; $c$ representa el contexto conservado para el modelo; $\theta$ son los parámetros. Aunque la supervisión "proviene de los propios datos", $g$ y $c$ aún deciden qué información se incentiva que el modelo conserve.

\lecturefigure{slide-017.jpg}{Aprendizaje auto-supervisado: construcción del objetivo de entrenamiento desde el interior de los datos originales}{CS25 V6 oficial deck, p.~17}

\teachervoice{00:07:10--00:08:13，los ponentes resumen la línea histórica como características diseñadas a mano, aprendizaje de representación supervisado y auto-supervisión. El docente enfatiza que el aprendizaje auto-supervisado reduce las etiquetas manuales, pero no elimina el sesgo inductivo traído por el objetivo.}

\begin{warningbox}{"Sin etiquetas manuales" no equivale a "sin suposiciones humanas"}
La proporción de ocultamiento, la dirección de predicción, las muestras negativas, el aumento de datos y la longitud de la ventana alteran la señal aprendible. Si el aumento destruye factores relevantes para la tarea, un modelo contrastivo será recompensado por ignorar información verdaderamente útil; si los datos del siguiente token están llenos de plantillas y repeticiones, el modelo priorizará aprender formas de alta frecuencia. El aprendizaje auto-supervisado es una forma de construir etiquetas, no una ley natural neutral al valor.
\end{warningbox}

Desde la perspectiva del cuello de botella de información, el objetivo decide qué información merece ser conservada. La clasificación de sentimientos solo requiere conservar las direcciones relacionadas con la etiqueta, y puede descartar activamente entidades, secuencia temporal y detalles gramaticales; el objetivo de reconstrucción necesita conservar más información local, pero podría codificar también una gran cantidad de ruido de píxeles; la predicción del siguiente token debe comprimir las leyes estadísticas útiles para futuros tokens, pero no recibe recompensa directa por estados del mundo real que no cambian la probabilidad del token. Por tanto, al evaluar una representación no basta mirar solo una puntuación downstream, sino que se deben revisar sondas lineales, transferencia con pocos ejemplos (\textit{few-shot}), cambio de distribución e intervenibilidad. Una representación que funcione bien en clasificación no significa que sea adecuada para planificación; un modelo generativo que pueda reconstruir detalles no implica que su latente ya haya abstraído objetos y causalidad.

\subsection{Casos lingüísticos y visuales: el objetivo de predicción determina la granularidad de la representación}

El caso lingüístico pasa de la clasificación de sentimientos a la predicción del siguiente token. Lo primero solo exige comprimir un texto en pocas etiquetas; lo segundo requiere predecir la distribución condicional en cada posición, por lo que la supervisión es más densa. Sea $x_{1:T}$ la secuencia de tokens; el modelo de lenguaje auto-regresivo optimiza
\[
\mathcal{L}_{\mathrm{NTP}}(\theta)
=-\sum_{t=1}^{T}\log p_{\theta}(x_t\mid x_{<t}).
\]
$x_t$ es el $t$-ésimo token, $x_{<t}$ es el contexto histórico. Cada posición contribuye a la pérdida, pero el objetivo sigue requiriendo solo predecir la secuencia observada, sin exigir directamente que el modelo proporcione fuentes factuales, ejecute herramientas o mantenga memoria a largo plazo.

\lecturefigure{slide-018.jpg}{Caso lingüístico: paso de etiquetas esparsas de sentimiento a predicción densa de secuencias}{CS25 V6 official deck, p.~18}

\lecturefigure{slide-020.jpg}{Estado completo del auto-supervisión lingüística: predecir el siguiente token bajo condición del contexto}{CS25 V6 official deck, p.~20}

\teachervoice{00:08:13--00:09:57，el docente contrasta la clasificación de sentimientos con el modelado de lenguaje: la clasificación proporciona solo una etiqueta de baja banda ancha, mientras que la predicción de secuencias genera supervisión en cada posición. Experiencia práctica: una supervisión más densa no implica que el objetivo cubra automáticamente la factualidad y la causalidad.}

El auto-supervisión visual suele elegir entre reconstrucción y aprendizaje contrastivo. El objetivo de reconstrucción exige que el decodificador restituya la entrada, lo que facilita conservar detalles de píxeles; el objetivo contrastivo acerca las vistas aumentadas del mismo objeto y aleja muestras distintas, enfatizando más la representación invariante. Una pérdida InfoNCE simplificada es
\[
\mathcal{L}_{\mathrm{NCE}}
=-\log\frac{\exp(\operatorname{sim}(z_i,z_i^+)/\tau)}
{\sum_j \exp(\operatorname{sim}(z_i,z_j)/\tau)}.
\]
$z_i$ es la representación ancla, $z_i^+$ la muestra positiva y $z_j$ las muestras candidatas; $\tau$ es la temperatura. Lo clave no es memorizar la fórmula, sino comprender que la definición de muestras positivas y negativas decide qué cambios se exige al modelo ignorar.

\lecturefigure{slide-021.jpg}{Caso visual: dos interfaces de supervisión entre reconstrucción con autoencoder y aprendizaje contrastivo}{CS25 V6 official deck, p.~21}

\lecturefigure{slide-022.jpg}{Autoencoder enmascarado: predecir contenido visual faltante tras ocultar parches}{CS25 V6 official deck, p.~22}

\begin{knowledgebox}{Lectura de la figura: límites de evidencia entre reconstrucción y contraste}
Primero observe el espacio de salida: los modelos de reconstrucción predicen detalles de entrada, mientras que los modelos contrastivos predicen relaciones entre muestras. Luego examine la invariancia: el segundo comprimirá activamente los cambios considerados "poco importantes" en el aumento. La figura respalda las diferencias entre dos objetivos de representación, pero no permite afirmar solo con un diagrama cual de ambos es generalmente mejor; tareas downstream, escala de datos, aumentos y capacidad del decodificador alterarán la conclusión.
\end{knowledgebox}

\subsection{Resumen de la sección}

La evolución de los paradigmas de aprendizaje se puede comprimir en "el desplazamiento interno continuo de la interfaz de supervisión": primero los humanos crean características, luego las redes aprenden representaciones y finalmente los propios datos generan objetivos densos. Pero cualquier objetivo solo recompensa patrones estadísticos específicos. Antes de entender el Transformer, hay que comprender que recibe tokens, optimiza predicciones condicionales y no obtiene automáticamente hechos, acciones ni coherencia a largo plazo.


\section{Embedding, RNN y Transformer: cómo fluye la información en secuencias}

El capítulo anterior explicó de dónde provienen las señales de supervisión; este entra en los mecanismos de cálculo. Los modelos de secuencia deben resolver tres cosas: colocar símbolos discretos en un espacio continuo, permitir que posiciones diferentes intercambien información y preservar el orden. El embedding resuelve la primera, las RNN lo hacen con estados recursivos para las otras dos, mientras que Transformer reorganiza el flujo de información usando atención y posición.

\subsection{De token a vector: representación estática vs. contextual}

Supongamos un vocabulario de tamaño $|V|$ y una dimensión oculta de $d$, donde la tabla de embedding es $E\in\mathbb{R}^{|V|\times d}$. El vector inicial del token id $x_i$ es
\[
e_i=E[x_i],\qquad e_i\in\mathbb{R}^{d}.
\]
La tokenización corta primero la cadena en palabras, subpalabras o bytes; el embedding luego mapea los ids discretos a vectores entrenables. El embedding estático usa el mismo vector para un mismo token, mientras que la representación contextual se actualiza en capas posteriores según el contexto, lo que permite separar polisemias.

\lecturefigure{slide-023.jpg}{Word embeddings: proximidad semántica, aritmética de vectores y limitaciones contextuales}{CS25 V6 oficial deck, p.~23}

\teachervoice{00:10:14--00:11:31，el docente explica primero la necesidad del embedding diciendo que el modelo solo procesa números, luego señala la limitación de usar un único vector para una misma palabra en diferentes contextos. Nota de clase: la contextualización ocurre durante las interacciones secuenciales posteriores, no se completa con el tokenizer.}

\begin{knowledgebox}{Primera vez: token, embedding y estado contextual}
El \textbf{token} es una unidad discreta de secuencia; el \textbf{embedding} es el vector inicial aprendible del token; el \textbf{contextual state} es el estado de información ya mezclado con contexto después de múltiples capas de atención/MLP. Los tres corresponden respectivamente a límites de corte, consulta de parámetros y cálculo dinámico, no deben confundirse.
\end{knowledgebox}

\subsection{RNN: comprimir el pasado en un estado recursivo}

La Recurrent Neural Network (RNN, red neuronal recurrente) actualiza el estado oculto paso a paso en el tiempo:
\[
h_t=\phi(W_x x_t+W_h h_{t-1}+b),\qquad
\hat{y}_t=W_o h_t.
\]
$x_t$ es la entrada actual, $h_{t-1}$ es la compresión histórica, $W_x,W_h,W_o$ son parámetros y $\phi$ es no lineal. El LSTM mitiga el desvanecimiento de gradientes mediante puertas, pero cada paso sigue dependiendo del anterior; el entrenamiento es difícilmente paralelo y la información a distancia larga también debe atravesar caminos largos.

\lecturefigure{slide-024.jpg}{RNN/LSTM antes de Transformer: actualización recursiva del estado paso a paso}{CS25 V6 oficial deck, p.~24}

\teachervoice{00:11:31--00:12:19，el docente no presenta la RNN como una "arquitectura fallida", sino que enfatiza su sesgo inductivo secuencial y dependencia paso a paso. Sirve para estados de flujo, pero resulta difícil aprovechar el entrenamiento paralelo a gran escala.}

\begin{warningbox}{El estado recursivo no es memoria infinita}
La dimensión de $h_t$ está fija; la información histórica debe comprimirse. Aunque el LSTM pueda mejorar la propagación de gradientes, esto no significa que cualquier detalle temprano se pueda preservar sin pérdidas. El olvido, la cobertura y la asignación de crédito en secuencias largas son problemas estructurales, no solo de ingeniería al aumentar el tamaño oculto.
\end{warningbox}

\subsection{Autoatención: reescribir las interacciones secuenciales como búsqueda blanda}

En la sección anterior, la RNN comprimía toda la historia en un único estado recursivo; aquí surge la pregunta: ¿podemos permitir que la posición actual acceda directamente a toda la secuencia sin que la información deba atravesar gradualmente $h_{t-1}, h_{t-2}, \ldots$? La respuesta de autoatención es construir una matriz direccionable por contenido. Al leer las siguientes fórmulas, observe primero cómo se genera el puntaje y luego cómo se pondera el valor; los pesos de softmax son solo rutas de la capa actual, no un repositorio permanente de conocimiento.

Dado una matriz de entrada $X\in\mathbb{R}^{n\times d}$, el modelo aprende
\[
Q=XW_Q,\qquad K=XW_K,\qquad V=XW_V,
\]
y calcula
\[
\operatorname{Attn}(Q,K,V)
=\operatorname{softmax}\!\left(\frac{QK^{\top}}{\sqrt{d_k}}\right)V.
\]
$Q$ representa la consulta de la posición actual, $K$ los índices disponibles para coincidir y $V$ el contenido realmente agregado; $d_k$ es la dimensión de la clave. Dividir por $\sqrt{d_k}$ controla la escala del producto punto para evitar que softmax se sature prematuramente.

\lecturefigure{slide-025.jpg}{Transformer en dos minutos: autoatención permite a cada token acceder a otras posiciones}{CS25 V6 oficial deck, p.~25}

\lecturefigure{slide-026.jpg}{Analogía de la biblioteca QKV: consulta, índice resumen y contenido del libro}{CS25 V6 oficial deck, p.~26}

\teachervoice{00:12:19--00:14:20，el docente usa una biblioteca para explicar QKV: query es lo que se busca, key es el resumen o índice del libro, y value es la información recuperada. Recordatorio de las apuntes: Q/K/V son proyecciones lineales diferentes de la misma entrada, no tres bases de datos independientes.}

\begin{lstlisting}[language=Python,caption={Pseudocódigo mínimo para atención de producto punto escalado}]
def attention(x, Wq, Wk, Wv):
    q = x @ Wq
    k = x @ Wk
    v = x @ Wv
    score = q @ k.T / sqrt(k.shape[-1])
    weight = softmax(score, axis=-1)
    return weight @ v
\end{lstlisting}

\subsection{Tokenización, posición y cabezas múltiples: las tres interfaces de un mismo bloque}

La atención por sí sola solo ve conjuntos de vectores. El modelo también necesita que el tokenizer decida la granularidad de la secuencia, que la codificación posicional inyecte el orden y que la estructura multicabeza permita establecer relaciones en paralelo en diferentes subespacios. Si el texto original se corta en $n$ tokens, la matriz de puntajes de autoatención estándar es $n\times n$, con un costo de cómputo y memoria aproximado que crece como $n^2$; esto es también una de las fuentes principales del costo en sistemas de contexto largo.

\lecturefigure{slide-027.jpg}{Tokenización: una misma frase puede cortarse en unidades secuenciales diferentes}{CS25 V6 oficial deck, p.~27}

\lecturefigure{slide-028.jpg}{Codificación posicional: inyectar orden en atención invariante a permutaciones}{CS25 V6 oficial deck, p.~28}

El vector de posición puede escribirse como $p_i$, con estado inicial $z_i^{(0)}=e_i+p_i$. La codificación sinusoidal clásica usa frecuencias distintas:
\[
p_{i,2k}=\sin\!\left(i/10000^{2k/d}\right),\qquad
p_{i,2k+1}=\cos\!\left(i/10000^{2k/d}\right).
\]
$i$ es la posición, $k$ el índice de frecuencia y $d$ la dimensión. Los modelos modernos también usan RoPE, ALiBi o mecanismos de posición relativa, pero el objetivo común es que los pesos de interacción perciban el orden relativo o absoluto.

La atención multicabeza divide los canales en múltiples cabezas:
\[
\operatorname{MHA}(X)=\operatorname{Concat}(H_1,\ldots,H_h)W_O,
\quad H_j=\operatorname{Attn}(XW_Q^{(j)},XW_K^{(j)},XW_V^{(j)}).
\]
$h$ es el número de cabezas; cada una tiene proyecciones independientes y $W_O$ se encarga de volver a mezclar. Las cabezas no corresponden automáticamente a relaciones gramaticales nombrables; simplemente proporcionan varios subespacios de interacción paralelos.

\lecturefigure{slide-029.jpg}{Atención multicabeza: múltiples conjuntos de relaciones de consulta se agregan en paralelo y luego se mezclan}{CS25 V6 oficial deck, p.~29}

\lecturefigure{slide-030.jpg}{Transformer completo: apilamiento de atención, MLP, residuales y normalización}{CS25 V6 oficial deck, p.~30}

\begin{knowledgebox}{Lectura de la figura: de la tabla de componentes al flujo de información}
Primero siga el flujo residual principal; luego observe cómo la atención intercambia información entre tokens y cómo el MLP transforma los canales posición por posición. LayerNorm estabiliza la escala. Cada bloque en la figura no es un "módulo de comprensión"; la capacidad proviene de la representación formada conjuntamente por muchas capas, datos de entrenamiento y objetivos. Visualizar una sola cabeza solo proporciona pistas sobre correlaciones locales, no constituye automáticamente una explicación causal.
\end{knowledgebox}

\subsection{Por qué Transformer reemplazó a RNN y por qué el problema no terminó}

Las ventajas de Transformer provienen de caminos cortos y cómputo paralelo: cualquier par de tokens puede interactuar en una sola capa de atención, y durante el entrenamiento se procesan todas las posiciones simultáneamente. Cada paso de la RNN tiene un costo menor y es naturalmente de flujo, pero las dependencias a larga distancia requieren múltiples recursiones. Si la dimensión oculta es $d$ y la longitud es $n$, la interacción de atención es aproximadamente $O(n^2d)$, mientras que las proyecciones posición por posición/MLP son aproximadamente $O(nd^2)$; el cuello de botella real depende de $n$, $d$, batch y hardware.

\lecturefigure{slide-031.jpg}{RNN vs. Transformer: contraste entre caminos recursivos, entrenamiento paralelo y dependencias largas}{CS25 V6 oficial deck, p.~31}

\lecturefigure{slide-032.jpg}{Mapa actual de Transformer: lenguaje, visión, voz, biología, video y robótica}{CS25 V6 oficial deck, p.~32}

\lecturefigure{slide-033.jpg}{Large Language Model: escala, preentrenamiento, autoregresivo y aprendizaje por contexto}{CS25 V6 oficial deck, p.~33}

\teachervoice{00:14:20--00:16:20，el docente atribuye el éxito al entrenamiento paralelo, modelado de dependencias largas e interfaz unificada de token, señalando al mismo tiempo que la atención cuadrática y el razonamiento de contexto largo siguen siendo limitaciones. Nota de clase: una arquitectura dominante no significa que esté en la frontera de Pareto óptima para todos los workloads.}

\begin{importantbox}{Explicación correcta de "Transformer está hoy en todas partes"}
La interfaz de secuencia unificada permite que texto, parches de imagen, marcos de audio, residuos de proteínas y tokens de acción compartan un backbone similar; además, la infraestructura de entrenamiento madura reduce el costo de replicación entre dominios. Sin embargo, cada modalidad aún requiere tokenizador/encoder especializado, receta de datos, pérdida, evaluación y restricciones de serving. Compartir bloques no significa que los problemas estén homogeneizados.
\end{importantbox}

Desde la contabilidad de recursos, las ventajas y costos de Transformer provienen de la misma cosa: establecer explícitamente interacciones token a token. Durante el entrenamiento se debe guardar o recalculan las activaciones; en la inferencia, los modelos autoregresivos usualmente guardan KV cache; el crecimiento de la secuencia aumenta simultáneamente la presión en cómputo de puntajes, capacidad del caché y ancho de banda de memoria. En secuencias cortas con dimensiones ocultas grandes, MLP y proyecciones pueden dominar los FLOPs principales; con secuencias ultra largas o muchos tokens multimodales, atención y KV cache dominan rápidamente. No basta en ingeniería decir "atención es $O(n^2)$"; hay que especificar batch, cabezas, tipo de dato, fases prefill/decode y ancho de banda del hardware. La razón por la cual RAG, agentes y modelos de video son costosos a menudo no es un aumento repentino de parámetros del modelo, sino que ponen más evidencia, historia y observaciones en el contexto.

Aquí también hay que distinguir entrenamiento paralelo de estado de inferencia. Durante el entrenamiento toda la secuencia es conocida, por lo que se puede construir $Q,K,V$ una vez y hacer multiplicación de matrices; durante el decode autoregresivo cada paso añade solo una consulta, pero debe leer todas las claves/valores anteriores. KV cache evita cálculos repetidos de tokens históricos, pero convierte parcialmente el problema de cómputo en uno de capacidad y ancho de banda de memoria. Multi-Query/Grouped-Query Attention reduce el caché permitiendo que múltiples cabezas de query compartan menos cabezas K/V; windowed o sparse attention reducen costos limitando posiciones accesibles; ambas pueden sacrificar acceso granular a largo plazo. El diseño del sistema es entonces un compromiso entre calidad, acceso aleatorio, ancho de banda y latencia, no una simple elección de "atención completa o ninguna".

\subsection{Resumen de la sección}

El embedding coloca tokens en espacio continuo, las RNN transmiten historia mediante estados recursivos y Transformer reescribe el flujo de información secuenciales con búsqueda blanda QKV, mecanismos de posición e interacción multicabeza. Su escalabilidad proviene de cálculo matricial paralelo e interfaz unificada; sus limitaciones provienen de interacciones cuadráticas, ejecución de contexto y restricciones de datos y sistema. El próximo capítulo gira hacia la otra mitad que determina el contenido del gradiente: estrategia de datos de preentrenamiento.

\section{Datos de preentrenamiento: más allá del volumen, origen, secuencia y memoria externa}

El capítulo anterior explicó cómo calcula el modelo; este pregunta de qué distribución provienen realmente los gradientes. Para la distribución de datos de entrenamiento $q(x)$, el modelo optimiza en la práctica
\[
\mathbb{E}_{x\sim q}\big[\ell_{\theta}(x)\big].
\]
Cambiar $q$ altera los hechos, lenguajes, estilos y patrones de razonamiento que el modelo ve repetidamente. El número de tokens es solo una dimensión; la calidad, la cobertura, la tasa de duplicación, el orden, las interfaces de recuperación y el schedule de crecimiento de parámetros también modifican la supervisión efectiva. La clase divide "los datos son importantes" en cuatro proyectos de investigación comprobables.

\subsection{Los datos no son un barril de petróleo, sino parte del objetivo de entrenamiento}

Llamar a los datos "nuevo petróleo" lleva a pensar solo en la cantidad. Una afirmación más precisa es: los datos definen la distribución de muestreo del riesgo empírico. Si en el corpus web cierto tipo de plantilla, hechos erróneos o proporción lingüística son excesivos, el optimizador los tratará como leyes estables; si las muestras de alto valor son escasas, incluso con suficiente capacidad del modelo, los gradientes no necesariamente reforzarán esa habilidad repetidamente. Por ello, la curaduría de datos, mixture y schedule son componentes externos de la arquitectura.

\lecturefigure{slide-035.jpg}{Impacto de los datos en IA: calidad, volumen, potencia computacional y capacidad acoplados}{CS25 V6 oficial deck, p.~35}

\lecturefigure{slide-036.jpg}{Proyecto de estrategia de datos uno: entrada individual infantil y aprendizaje bilingüe}{CS25 V6 oficial deck, p.~36}

\lecturefigure{slide-037.jpg}{Proyecto de estrategia de datos dos: asignación computacional RAG y escalado guiado por currículo}{CS25 V6 oficial deck, p.~37}

\teachervoice{00:16:20--00:20:23，el docente divide la investigación de datos en cuatro tipos de intervención: cambiar el origen, cambiar la mezcla lingüística, externalizar el conocimiento a recuperación externa y hacer variar conjuntamente la dificultad de los datos con la profundidad del modelo. La diapositiva recuerda: estos proyectos responden preguntas causales distintas y no pueden sintetizarse en una frase "mejores datos, mejor modelo".}

\begin{knowledgebox}{Lectura de la figura: cuatro proyectos que modifican diferentes variables}
Baby Scale modifica la \textbf{fuente de datos familiar/infantil}; Bilingual BabyLM modifica la \textbf{estructura de mezcla lingüística y hablantes}; RAG scaling modifica la \textbf{asignación computacional entre memoria en parámetros y memoria externa}; CGLS modifica el \textbf{orden temporal entre profundidad del modelo y dificultad de las muestras}. Primero se define la intervención, luego se observa el métrica y solo después se pueden discutir conclusiones transferibles.
\end{knowledgebox}

\begin{importantbox}{Los cinco perillas de la estrategia de datos}
\begin{enumerate}
  \item \textbf{Selection}: qué muestras entran al corpus;
  \item \textbf{Mixture}: cómo se proporcionalizan las lenguas, dominios y tareas;
  \item \textbf{Order}: cuándo aparecen los datos simples/difíciles o antiguos/nuevos;
  \item \textbf{Externalization}: si el conocimiento va a los parámetros o al store de recuperación;
  \item \textbf{Co-scaling}: si el schedule de datos varía sincronizado con parámetros, capas y potencia computacional.
\end{enumerate}
Estas perillas modifican los gradientes y no deben verse como una limpieza única antes del entrenamiento.
\end{importantbox}

Estos cuatro proyectos también ilustran el diseño causal básico de la investigación de datos. Si se cambia el origen de los datos, se debe fijar el tokenizer, el número de parámetros, los tokens de entrenamiento y el optimizador; si se cambia la mezcla lingüística, se debe reportar simultáneamente la distribución de prueba por cada idioma; si se cambia el presupuesto de recuperación, se deben incluir el contexto adicional y los costos de índice; si se cambia el schedule de capas, se debe controlar la cantidad total de actualización de parámetros y la potencia computacional de entrenamiento. De lo contrario, el llamado "efecto de datos" podría ser solo que el modelo es más grande, el entrenamiento más largo o la evaluación más cercana a la distribución de entrenamiento. La clase no escribió estos experimentos como una ley universal, sino que proporcionó un método: definir explícitamente un conjunto de intervenciones por vez y escribir los confound sin controlar en los límites de la conclusión.

\subsection{Resumen de la sección}

Los datos de preentrenamiento determinan no solo qué "ve" el modelo, sino también la frecuencia con que se ven las muestras, su orden, contexto y recuperabilidad. Los cuatro casos siguientes responderán por separado: qué aprenden los datos pequeños y naturales, si el bilingüe se interfiere mutuamente, cuánto presupuesto debe tomar RAG y si el crecimiento del modelo puede sincronizarse con el currículo.
\section{Datos pequeños y aprendizaje bilingüe: convertir la intuición humana en experimentos comprobables}

La clase utiliza el aprendizaje infantil para plantear la cuestión de la eficiencia con datos, pero no equipara las redes neuronales con el cerebro. Los humanos reciben entrada visual, motora, interacción social, metas a largo plazo y alta estructuración; los modelos de lenguaje suelen recibir tokens offline y una función de pérdida unificada. El valor de la investigación reside en construir condiciones de datos controladas, no en demostrar directamente con pocos tokens un "aprendizaje tipo cerebro".

\subsection{¿Por qué estudiar small model y small data?}

Baby Scale compara infant/ChatGPT con brain/neural network, listando aprendizaje continuo, input grounded, metas, estructura jerárquica y sesgo inductivo latente. Esta comparación actúa como generador de hipótesis: si la interactividad, la diversidad semántica o la estructura del hablante en la entrada infantil difieren de la de las páginas web, se puede diseñar una prueba a nivel de dataset; pero no se deben convertir todas las supervisiones humanas a caracteres de transcripción de texto.

\lecturefigure{slide-039.jpg}{Infantes versus ChatGPT, cerebro versus red neuronal: contraste entre dos sistemas de aprendizaje}{CS25 V6 official deck, p.~39}

\lecturefigure{slide-040.jpg}{Diferencias en continuidad, metas, modalidades y estructura entre el aprendizaje humano y el entrenamiento de LLMs}{CS25 V6 official deck, p.~40}

\lecturefigure{slide-041.jpg}{Motivación para estudiar modelos pequeños y datos pequeños: eficiencia, interpretabilidad y experimentos controlados}{CS25 V6 official deck, p.~41}

\teachervoice{00:18:22--00:20:46，el docente utiliza repetidamente might、potential y question para describir las ventajas del aprendizaje humano. Nota de clase: estas diferencias son ejes experimentales, no la conclusión de que "los datos infantiles son naturalmente superiores a los de internet".}

\begin{warningbox}{¿Qué falta al comparar el número de tokens entre infantes y LLMs?}
Los infantes reciben simultáneamente imágenes, sonido, consecuencias de acciones, retroalimentación social e interacción repetida; el corpus transcrita conserva solo una parte. Las funciones objetivo, la inicialización de parámetros, los conocimientos previos y las tareas de evaluación también difieren entre modelos e infantes. Por tanto, el conteo de tokens puede plantear la cuestión de la eficiencia con datos, pero no representa una cantidad total de supervisión justa.
\end{warningbox}

\subsection{Baby Scale: ¿qué puede entrenar la entrada lingüística de un solo niño?}

Baby Scale utiliza speech dirigida al niño (child-directed speech) de una sola familia del corpus CHILDES, entrena modelos de estilo GPT-2 pequeños y compara en tareas de language modeling, sintaxis, semántica y desarrollo lexical. El diseño central no persigue SOTA, sino observar si el \textbf{origen familiar de los datos} y la \textbf{estructura de exposición infantil} generan diferencias medibles una vez fija la escala del modelo.

\lecturefigure{slide-043.jpg}{Pregunta central de Baby Scale: ¿puede un modelo aprender de la entrada de un solo niño?}{CS25 V6 official deck, p.~43}

\lecturefigure{slide-044.jpg}{Configuración experimental: transcripciones familiares, escala de tokens, modelos y tareas de evaluación}{CS25 V6 official deck, p.~44}

La perplexity es la exponencial de la entropía cruzada:
\[
\operatorname{PPL}=\exp\!\left(-\frac{1}{T}\sum_{t=1}^{T}\log p_{\theta}(x_t\mid x_{<t})\right).
\]
Intuitivamente, aproxima el número "efectivo de candidatos" frente a cada token; valores más bajos indican menos sorpresa ante la secuencia de prueba. Sin embargo, PPL está influenciado por el tokenizer, la distribución de prueba y la longitud; no permite comparación directa entre modelos arbitrarios ni equivale a capacidad gramatical, factual o interactiva.

La página de resultados debe leerse en dos pasos: la diapositiva 45 compara el rendimiento relativo de diferentes datos familiares en múltiples tareas, mientras que la 46 revisa si el tamaño del modelo altera la tendencia. Si un conjunto de datos es fuerte en varias tareas, puede deberse a cobertura lexical, estructura interactiva o diversidad del hablante; si ampliar el modelo no elimina las diferencias, indica que las propiedades de los datos no son simplemente un cuello de botella de capacidad.

\lecturefigure{slide-045.jpg}{Resultados multitarea de Baby Scale: diferentes entradas familiares generan perfiles de capacidades distintos}{CS25 V6 official deck, p.~45}

\lecturefigure{slide-046.jpg}{Interacción entre escala del modelo y origen de datos: aumentar parámetros no uniformiza ni elimina las diferencias}{CS25 V6 official deck, p.~46}

\begin{knowledgebox}{Lectura de la figura: comparar primero tareas, luego escalas}
En primer lugar, compare diferentes datos familiares dentro de cada subplot, en lugar de comparar puntuaciones absolutas entre tareas; los métricas y niveles de dificultad varían por tarea. Luego, observe si el ordenamiento se mantiene estable tras cambiar el tamaño del modelo. Un ordenamiento estable apoya la idea de que "las propiedades de los datos afectan continuamente al rendimiento", pero no basta con la correlación para determinar si es la interacción infantil, la cobertura lexical o la calidad de transcripción la causa; se requieren ablaciones controladas posteriores.
\end{knowledgebox}

\subsection{Más allá de la escala: ¿qué propiedades de los datos predicen el rendimiento?}

La clase lista candidatos como interaction, diversidad del hablante, longitud de contexto, cobertura lexical, proporción child/caregiver speech y divergencia de distribución. Esto puede expresarse desde una perspectiva de regresión simple:
\[
y_{f,m}=\beta_0+\beta_1\log D_f+\beta_2\log N_m
+\sum_k\gamma_k z_{f,k}+\varepsilon_{f,m},
\]
donde $y_{f,m}$ es el rendimiento de la tarea para la familia $f$ y escala del modelo $m$, $D_f$ es el número de tokens, $N_m$ son los parámetros y $z_{f,k}$ representa las propiedades de los datos. Esta ecuación no es el único modelo del artículo, sino una ayuda para distinguir efectos de tamaño (size effect) de efectos de composición (composition effect).

\lecturefigure{slide-047.jpg}{Resultados de Scaling: el impacto del dataset y la escala del modelo en tareas distintas no es uniforme}{CS25 V6 official deck, p.~47}

\lecturefigure{slide-048.jpg}{Más allá del tamaño: variables explicativas candidatas como interacción, vocabulario, hablante y contexto}{CS25 V6 official deck, p.~48}

\lecturefigure{slide-049.jpg}{Conclusión general: los datos familiares con mejor rendimiento son más ricos, interactivos y estructuralmente semánticos}{CS25 V6 official deck, p.~49}

\lecturefigure{slide-050.jpg}{Takeaways de Baby Scale: el aprendizaje a escala infantil es viable pero altamente dependiente de la tarea}{CS25 V6 official deck, p.~50}

\teachervoice{00:22:23--00:25:25，el docente enfatiza que "más grande" no es la única explicación: la proporción child/caregiver, la diversidad interactiva y la estructura lexical pueden afectar el rendimiento; además, los datasets óptimos varían según la tarea.}

\begin{importantbox}{Conclusiones que Baby Scale sí y no apoya}
Apoya que: modelos pequeños pueden aprender patrones lingüísticos no aleatorios a partir de transcripciones de un entorno infantil individual; las diferencias en el origen de los datos dejan huellas relevantes para la tarea. No apoya que: la entrada infantil explique por completo el aprendizaje humano del lenguaje; cualquier corpus familiar sea óptimo necesariamente en modelos grandes; o que las propiedades de los datos correlacionadas se hayan demostrado como mecanismos causales.
\end{importantbox}

Al evaluar este tipo de investigación, también hay que evitar "un único puntaje que cubra todas las capacidades". La función de pérdida de language modeling mide predicciones locales; tareas de grammaticality del tipo BLiMP verifican preferencias estructurales; la similitud lexical revisa la representación lexical; y las medidas de desarrollo se interesan por la trayectoria a medida que crece el dato. Un corpus familiar puede dominar en tareas semánticas pero rezagarse en sintácticas; esto no es una contradicción, sino que indica que las propiedades de los datos soportan capacidades distintas. Experimentos posteriores más fuertes deberían realizar emparejamiento o intervención sobre propiedades altamente correlacionadas, por ejemplo re-muestreando diversidad del hablante manteniendo el número de tokens fijo, desordenando turnos de interacción o controlando cobertura lexical, para ver si los perfiles de tareas cambian según lo predicho. Solo así se puede avanzar de "¿qué dataset es mejor" a "¿qué mecanismo de datos está actuando".

El tamaño de la muestra y el entrenamiento repetido también deben ingresar en la incertidumbre. Los experimentos con modelos pequeños suelen verse afectados por random seed, orden del batch y early stopping; si solo se reporta una ejecución, el ordenamiento del dataset podría ser ruido de optimización. Un reporte ideal debería dar medias multi-seed, intervalos de confianza, tamaños de efecto (effect size) y corrección para comparaciones múltiples, además de explicar si existen superposiciones o diferencias de calidad en las transcripciones a nivel familiar. Para las afirmaciones de desarrollo, también se debe comparar la curva de crecimiento del dato en lugar de mirar solo el checkpoint final. Solo así se puede juzgar si un corpus es consistentemente más efectivo o simplemente converge más rápido bajo un presupuesto de entrenamiento limitado.

\subsection{Bilingual BabyLM: ¿es el bilingüismo interferencia, transferencia o problema de escala de datos?}

La investigación bilingüe desglosa aún más "mezclar dos idiomas" en estructura de exposición. Se puede permitir que cada hablante use solo un idioma (speaker-split), que el hablante cambie entre oraciones (sentence-level) o dentro de la misma oración (word-level code-switching). Si todos los casos se llaman bilingual mixture, no se puede determinar a qué es sensible el modelo: límites lingüísticos, pistas del hablante o volumen total de datos.

\lecturefigure{slide-052.jpg}{Tres preguntas de investigación de Bilingual BabyLM: interferencia, estructura de exposición y escala}{CS25 V6 official deck, p.~52}

\lecturefigure{slide-053.jpg}{Condiciones de entrenamiento: monolingüe, bilingüe aleatorio, cambio en nivel speaker/sentence/word}{CS25 V6 official deck, p.~53}

Sea $q_{\mathrm{en}}$ y $q_{\mathrm{es}}$ la distribución de datos en inglés y español, con peso de mezcla $\lambda$; entonces la distribución de entrenamiento es
\[
q_{\lambda}(x)=\lambda q_{\mathrm{en}}(x)+(1-\lambda)q_{\mathrm{es}}(x).
\]
Sin embargo, speaker-split y code-switching no solo cambian $\lambda$; alteran la distribución conjunta de lenguaje/contexto, hablante y límites de oración. Incluso con el mismo total de tokens, las pistas de límite disponibles para el modelo difieren.

\subsection{Resultado uno: no hay interferencia bilingüe universal}

El diseño experimental anterior desglosa la exposición bilingüe en múltiples condiciones; esta sección responde primero a la pregunta más básica: ¿aprender simultáneamente dos idiomas destruye sistemáticamente las capacidades monolingües? Al leer las diapositivas 55--57, primero se debe fijar el total de datos y el baseline, luego verificar por separado language modeling, sintaxis y semántica; solo si múltiples métricas degeneran significativamente en la misma dirección se puede apoyar una interferencia amplia.

La tabla de PPL en la slide 55 muestra que los modelos multilingües pueden formar distribuciones efectivas en ambos idiomas; las slides 56--57 luego usan grammaticality y word similarity para verificar si solo recuerdan estadísticas locales. Al leer la tabla, primero busque el baseline monolingüe, luego revise los intervalos de confianza de las condiciones bilingües, no solo compare una cifra decimal.

\lecturefigure{slide-055.jpg}{Perplexity de modelos bilingües en inglés y español: pueden modelar ambos idiomas simultáneamente}{CS25 V6 official deck, p.~55}

\lecturefigure{slide-056.jpg}{Grammaticality: comparación de capacidad estructural entre baseline en inglés y condiciones multilingües}{CS25 V6 official deck, p.~56}

\lecturefigure{slide-057.jpg}{Word similarity: contraste de representación semántica entre entrenamiento monolingüe y multilingüe}{CS25 V6 official deck, p.~57}

\begin{knowledgebox}{Lectura de la figura: cómo juzgar interferencia}
La interferencia no debe definirse solo como una condición ligeramente inferior al máximo puntaje. Se deben revisar simultáneamente el baseline monolingüe, el volumen total de datos, los intervalos de confianza y múltiples tareas. Si un modelo bilingüe baja ligeramente en inglés pero aprende español, el tradeoff general difiere de la "interferencia catastrófica"; si tras fijar el total de datos aún hay una caída significativa, esto se acerca más a problemas de capacidad competitiva u optimización.
\end{knowledgebox}

\subsection{Resultado dos: la estructura de exposición tiene un impacto menor al esperado}

Las curvas de grammaticality y word similarity para diferentes condiciones multilingües están globalmente cercanas, indicando que los modelos no dependen fuertemente de pistas humanas comunes como "un hablante por idioma". Este resultado debe interpretarse así: bajo el modelo actual, volumen de datos y métricas, el impacto marginal de la estructura de exposición es limitado; no se puede concluir que code-switching sea irrelevante en todos los idiomas, escalas o tareas generativas.

\lecturefigure{slide-059.jpg}{Rendimiento de grammaticality bajo diferentes estructuras de exposición bilingüe}{CS25 V6 official deck, p.~59}

\lecturefigure{slide-060.jpg}{Rendimiento de word similarity bajo diferentes estructuras de exposición bilingüe}{CS25 V6 official deck, p.~60}

\subsection{Resultado tres: la escala total de datos es más crítica que una pequeña expansión del modelo}

La slide 62 compara por separado aproximadamente 1/20 de la escala de parámetros y 1/20 de la escala de datos. Los resultados enfatizan el rol de data scale: reducir datos degrada significativamente múltiples métricas, mientras que en este rango de modelos pequeños, modest model-size change tiene menor impacto. Aquí aún hay que evitar extrapolar a LLMs frontier; en escalas mayores, la proporción óptima entre parámetros, datos y compute podría ser distinta.

\lecturefigure{slide-062.jpg}{Contraste entre escala de datos y modelo: reducir datos causa degradación más evidente}{CS25 V6 official deck, p.~62}

\lecturefigure{slide-063.jpg}{Conclusión RQ1: no se detectó interferencia bilingüe universal}{CS25 V6 official deck, p.~63}

\lecturefigure{slide-064.jpg}{Conclusión RQ2: la estructura de speaker y code-switching tiene impacto menor al esperado}{CS25 V6 official deck, p.~64}

\lecturefigure{slide-065.jpg}{Conclusión RQ3: en el rango actual, data scale es más importante que modest model size}{CS25 V6 official deck, p.~65}

\lecturefigure{slide-066.jpg}{Takeaways de Bilingual BabyLM y límites experimentales}{CS25 V6 official deck, p.~66}

\teachervoice{00:25:25--00:29:01，el docente resume tres puntos: no hay interferencia monolingüe amplia; el impacto de la estructura de exposición es sorpresivamente débil; y la escala total de datos es más crítica que el tamaño del modelo en este rango. El docente simultáneamente limita las conclusiones a un small-scale setting.}

\begin{warningbox}{“El bilingüismo no interfiere” no es una ley condicional}
La distancia lingüística, el tokenizer, la proporción de mezcla, la duración del entrenamiento, las tareas generativas y la capacidad del modelo pueden cambiar los resultados. Lo que muestran los experimentos actuales es que bajo un conjunto controlado de condiciones small-scale no hay degeneración sistemática, no que todo multilingual training genere automáticamente transferencia positiva.
\end{warningbox}

Para la práctica de ingeniería, este conjunto de resultados se asemeja más a un orden de depuración que a una receta de mezcla. Si un modelo multilingüe degrada, primero revise el número efectivo de tokens por idioma, fertilidad del tokenizer, muestreo del batch y filtrado de evaluación (evaluation leakage), luego verifique competencia de capacidad. Speaker-split o estructura de code-switching pueden volverse importantes en tareas contextuales más fuertes; las métricas actuales de sintaxis y similitud lexical no son sensibles al cambio lingüístico a nivel de discurso. La prudencia del curso radica en separar "no se observó un efecto significativo" de "el efecto no existe", lo cual es el principio básico para leer cualquier resultado negativo.

\subsection{Resumen de la sección}

Ambos proyectos BabyLM demuestran conjuntamente que el valor de los experimentos small-scale reside en separar variables de datos. Baby Scale muestra que las propiedades de los datos familiares se correlacionan con el rendimiento de la tarea; Bilingual BabyLM indica que la mezcla lingüística no necesariamente causa interferencia y que la estructura de exposición y el volumen total de datos juegan roles distintos. Proporcionan pistas sobre mecanismos, no una fórmula universal para replicar directamente el aprendizaje infantil en modelos grandes.
\section{RAG y Curriculum: integrar la ubicación del conocimiento y el orden de aprendizaje en el scaling}

El capítulo anterior modificó el origen y la estructura lingüística del corpus; este capítulo cambia dos dimensiones adicionales: si el conocimiento se almacena en los parámetros o se recupera durante la inferencia, y si el modelo crece junto con un curriculum y un schedule de profundidad. Ambas rutas se oponen a la visión unidimensional de scaling basada únicamente en "ampliar el preentrenamiento estático".

\subsection{La interfaz básica de RAG: agregar memoria externa más allá de la memoria paramétrica}

Retrieval-Augmented Generation (RAG, generación aumentada con recuperación) recupera primero evidencias del repositorio de documentos según la consulta y luego entrega tanto la consulta como las evidencias al modelo de lenguaje. Desplaza parte de la actualización del conocimiento desde el entrenamiento paramétrico hacia la actualización de índices, pero incrementa la latencia de recuperación, los tokens de contexto, los errores de clasificación y las contradicciones en las evidencias.

\lecturefigure{slide-068.jpg}{Modelo base: una vez que los datos de preentrenamiento se comprimen en parámetros, generan la salida directamente}{CS25 V6 oficial deck, p.~68}

\lecturefigure{slide-070.jpg}{RAG completo: la consulta activa la recuperación y las evidencias entran al LLM como contexto}{CS25 V6 oficial deck, p.~70}

Si el presupuesto total de cómputo es $C$, se puede expresar aproximadamente como
\[
C=C_{\mathrm{pretrain}}+C_{\mathrm{index}}+C_{\mathrm{retrieve}}+C_{\mathrm{generate}}.
\]
$C_{\mathrm{pretrain}}$ forma el conocimiento paramétrico, $C_{\mathrm{index}}$ construye la base de datos externa, $C_{\mathrm{retrieve}}$ realiza búsquedas en cada solicitud y $C_{\mathrm{generate}}$ procesa el contexto aumentado. RAG no es conocimiento gratuito; redistribuye los costos desde un entrenamiento único hacia el almacenamiento y cada paso de inferencia.

\begin{knowledgebox}{Uso inicial: memoria paramétrica y no paramétrica}
La \textbf{memoria paramétrica} se refiere a las leyes estadísticas escritas en los pesos mediante gradientes; su actualización es lenta, la consulta es barata pero el origen no se puede localizar directamente. La \textbf{memoria no paramétrica} se refiere a repositorios de documentos, índices vectoriales o bases de datos; permite actualizaciones rápidas y rastreo del origen, pero depende de la calidad de la recuperación, los índices y el contexto adicional. RAG es una combinación de ambas, no "aprendizaje sin parámetros".
\end{knowledgebox}

\subsection{Scaling de RAG: el óptimo depende de la tarea y del presupuesto}

La investigación integra modelos base como OLMo-2 con módulos de recuperación en una asignación unificada de cómputo. Si hay demasiado poco preentrenamiento, el modelo podría no entender la consulta, las evidencias o la tarea de generación; si hay demasiado preentrenamiento, los beneficios marginales de continuar presionando hechos recuperables hacia los parámetros disminuyen. La proporción óptima de recuperación varía según la tarea: las problemas intensivos en conocimiento dependen más de las evidencias externas, mientras que el modelado lingüístico puro o la inferencia compleja siguen requiriendo un modelo base fuerte.

\lecturefigure{slide-071.jpg}{Estudiando el Scaling de RAG: escala paramétrica, tokens de preentrenamiento y recuperación cambian conjuntamente}{CS25 V6 oficial deck, p.~71}

\lecturefigure{slide-072.jpg}{¿Cuánto RAG es óptimo?: la asignación óptima de recuperación difiere según la tarea}{CS25 V6 oficial deck, p.~72}

\teachervoice{00:29:01--00:31:50, el docente enfatiza que existe un presupuesto mínimo útil de preentrenamiento; RAG no puede hacer que un modelo casi sin entrenamiento se vuelva fuerte automáticamente mediante la recuperación. El óptimo depende de la tarea, el recuperador y el cómputo total, no de una proporción fija.}

\begin{lstlisting}[language=Python,caption={Pseudocódigo para explorar la asignación entre preentrenamiento y recuperación bajo presupuesto fijo}]
for pretrain_budget in candidate_pretrain_budgets:
    model = train_base_model(tokens=budget_to_tokens(pretrain_budget))
    for retrieval_k in candidate_k:
        score = evaluate(
            model=model,
            retriever=index,
            k=retrieval_k,
            task=target_task,
        )
        record(pretrain_budget, retrieval_k, score)
choose_pareto_frontier()
\end{lstlisting}

\begin{warningbox}{Los tres costos de RAG más fáciles de ignorar}
En primer lugar, recuperar documentos aumentará los tokens de entrada y el caché KV; en segundo lugar, las evidencias erróneas empujarán al modelo hacia respuestas incorrectas pero "fundadas"; en tercer lugar, la actualización del índice y el control de acceso se convierten en nuevos estados del sistema. Comparar únicamente el número de parámetros sin contar los costos de recuperación/servicio llevará a malinterpretar el óptimo.
\end{warningbox}

En el lado de servicio también hay que distinguir entre prefill y decode. Las $k$ documentos recuperados se expanden una sola vez la longitud de entrada durante la fase de prefill, aumentando los costos de atención y establecimiento del caché KV; después, cada token de decode sigue leyendo estos cachés. Si las solicitudes de usuario son de alta concurrencia, un contexto de evidencia más grande reducirá el batch soportable. Por otro lado, un modelo base débil podría no utilizar los documentos correctamente, produciendo copiado de citas, contradicciones en las evidencias o ignorando la recuperación. Por ello, un experimento completo de scaling de RAG debe reportar al menos recall de recuperación, calidad de respuesta, longitud de contexto, latencia, throughput y costo de actualización del índice fresco, no solo la precisión final.

\subsection{Curriculum-Guided Layer Scaling: modelo y datos crecen juntos}

El aprendizaje por curriculum ordena las muestras según cierta dificultad o estructura. CGLS lleva esto más allá permitiendo que el conteo de capas crezca con el progreso del entrenamiento: el modelo primero procesa datos más simples con una red más superficial, e inserta gradualmente capas y aumenta la complejidad de los datos. Sea $s\in[0,1]$ el progreso de entrenamiento; el schedule de capas y dificultad se puede expresar como
\[
L(s)=L_0+\lfloor s(L_{\max}-L_0)\rfloor,
\qquad d(s)=d_0+s(d_{\max}-d_0).
\]
$L(s)$ es la profundidad actual y $d(s)$ es la métrica de dificultad de los datos. Los métodos prácticos necesitan definir cómo replicar/inicializar nuevas capas, cómo mantener el estado del optimizador y si la dificultad es realmente medible.

\lecturefigure{slide-075.jpg}{Motivación de CGLS: un modelo estático podría no coincidir con la evolución de la dificultad de los datos durante el entrenamiento}{CS25 V6 oficial deck, p.~75}

\lecturefigure{slide-076.jpg}{Motivación: cursos sintéticos, desarrollo infantil y DataComp-LM proporcionan pistas de orden}{CS25 V6 oficial deck, p.~76}

\lecturefigure{slide-077.jpg}{Trabajo previo: métodos de apilado de capas y crecimiento impulsado por similitud}{CS25 V6 oficial deck, p.~77}

\lecturefigure{slide-078.jpg}{Método CGLS: crecimiento de 8 a 16 capas según la fase de entrenamiento con cambios en el curriculum}{CS25 V6 oficial deck, p.~78}

\begin{knowledgebox}{Lectura de la figura: no interpretar "curso" solo como ordenamiento de muestras}
La diapositiva 78 cambia simultáneamente la profundidad del modelo, las capas transferidas, las capas inicializadas aleatoriamente y la complejidad de los datos. Antes de leer los resultados, pregunte qué parámetros se heredan, replican o crean en cada fase; de lo contrario, los cambios de rendimiento podrían provenir del crecimiento paramétrico, pasos de entrenamiento o inicialización, no solo del curriculum en sí.
\end{knowledgebox}

\begin{lstlisting}[language=Python,caption={Pseudocódigo para el schedule conjunto de profundidad del modelo y dificultad de datos}]
for stage in schedule:
    model = grow_layers(
        model,
        target_depth=stage.depth,
        transfer=stage.transfer_policy,
    )
    data = sample_by_difficulty(
        corpus,
        min_level=stage.min_difficulty,
        max_level=stage.max_difficulty,
    )
    train(model, data, steps=stage.steps)
\end{lstlisting}

\subsection{Resultados y límites: los pequeños beneficios requieren atribución correcta}

Las diapositivas 79--80 resumen múltiples benchmarks. Al comparar, primero identifique una línea base fija de cómputo/parámetros y luego determine si CGLS lidera consistentemente en el promedio, perplexity o tareas específicas. Incluso si existen beneficios promedio, verifique si sacrifica tareas individuales y si la complejidad de implementación adicional del proceso de crecimiento vale la pena.

\lecturefigure{slide-079.jpg}{Tabla de resultados de CGLS I: cambios promedio y por tarea en múltiples benchmarks}{CS25 V6 oficial deck, p.~79}

\lecturefigure{slide-080.jpg}{Tabla de resultados de CGLS II: comparación con diferentes líneas base y configuraciones}{CS25 V6 oficial deck, p.~80}

\lecturefigure{slide-081.jpg}{Discusión: beneficios, línea base normalizada y direcciones futuras del curriculum}{CS25 V6 oficial deck, p.~81}

\teachervoice{00:31:50--00:36:33, el docente describe CGLS como un método exploratorio: los resultados mejoran, pero aún se requiere investigación en la definición del curriculum, hiperparámetros, otras familias de modelos y escalas mayores. Experiencia práctica: no eleve una tabla de promedios a "leyes de crecimiento humanas".}

El crecimiento de capas también introduce problemas de estado del optimizador. Adam/AdamW mantienen un primer momento $m$ y un segundo momento $v$ para cada parámetro; cuando se agregan nuevas capas, estos estados del optimizador no tienen historia. Copiar capas existentes puede heredar representaciones, pero podría causar homogeneización entre capas; la inicialización aleatoria mantiene la diversidad, pero crea saltos distribucionales durante el entrenamiento. Las comparaciones justas deben explicar cómo se inicializan los pesos de las nuevas capas, $m/v$, el warmup de la tasa de aprendizaje y la escala residual, y emparejar la línea base según FLOPs totales o cantidad total de actualizaciones. De lo contrario, los beneficios de CGLS podrían provenir de etapas adicionales de entrenamiento en lugar del crecimiento conjunto del curriculum y la profundidad.

\subsection{Perspectiva unificada de datos: selección, mezcla, recuperación y curriculum}

Los cuatro proyectos se pueden resumir en una tabla de decisión: Baby Scale estudia la selección de fuentes, Bilingual BabyLM estudia la estructura de mezcla, RAG estudia la externalización del conocimiento y CGLS estudia el orden y el co-scaling. Juntos demuestran que la estrategia de datos no es solo filtrar antes del entrenamiento, sino una variable sistémica que atraviesa el aprendizaje paramétrico, la inferencia y el crecimiento del modelo.

\lecturefigure{slide-084.jpg}{Insight unificado I: origen de los datos, mezcla lingüística y dependencia de tareas}{CS25 V6 oficial deck, p.~84}

\lecturefigure{slide-085.jpg}{Insight unificado II: ubicaciones de intervención diferentes para curriculum, RAG y crecimiento del modelo}{CS25 V6 oficial deck, p.~85}

\lecturefigure{slide-086.jpg}{Conclusión general: los futuros modelos de lenguaje necesitan ser más inteligentes, no solo tener más datos}{CS25 V6 oficial deck, p.~86}

\begin{importantbox}{Lista de selección de estrategias de datos}
Si el problema es \textbf{ruido en el corpus de entrenamiento}, haga primero selección/filtrado; si el problema es \textbf{desequilibrio lingüístico o de dominio}, ajuste la mezcla; si el conocimiento se actualiza frecuentemente y requiere fuentes, considere RAG; si las etapas tempranas y tardías del entrenamiento requieren diferentes muestras, diseñe un curriculum; si la capacidad del modelo cambia según la fase, debe incluir el crecimiento de arquitectura en el experimento. Los diferentes controles no pueden reemplazarse con una sola métrica de "calidad de datos".
\end{importantbox}

\subsection{Resumen de la sección}

RAG extiende la ubicación del conocimiento desde los parámetros hacia la memoria externa, y CGLS extiende el orden de aprendizaje hacia un schedule conjunto de profundidad del modelo y dificultad de los datos. Ambos revelan que el núcleo del scaling es la asignación de recursos, no aumentar incondicionalmente tokens o parámetros. El siguiente capítulo gira hacia post-entrenamiento e inferencia de cómputo en tiempo real que pueden cambiar el comportamiento después de completar el entrenamiento del modelo.
\section{Post-training y compute en tiempo de inferencia: cómo el comportamiento continúa cambiando después del preentrenamiento}

El \texttt{pretraining} permite al modelo adquirir una distribución condicional amplia, pero no convierte automáticamente las salidas en un asistente controlable. El \texttt{post-training} puede modificar los parámetros o añadir durante la inferencia prompts, búsqueda, recuperación, herramientas o verificadores. Al evaluar un método, pregúntese primero: ¿actualiza pesos, modifica el decodificación, añade estado externo o redefine las recompensas? De lo contrario, CoT, RLHF, RAG y agent corren el riesgo de mezclarse en una capa de "mejora" difusa.

\subsection{Mapa del post-training: ¿qué más hay que ajustar después del entrenamiento}

El curso coloca \texttt{post-training} y \texttt{inference-time compute} en la misma página: \texttt{prompt engineering}, ajuste de instrucciones (\texttt{instruction tuning}), optimización de preferencias, recuperación, agentes/herramientas pueden todos cambiar el comportamiento desplegado. Si la distribución del modelo base es $p_0(y\mid x)$, el post-training produce $p_{\theta}(y\mid x)$, y la estrategia en tiempo de inferencia puede escribirse como
\[
y^{\star}=\arg\max_{y\in\mathcal{S}(x;B)}
\left[s_{\theta}(y\mid x)+\lambda r(y,x)\right],
\]
donde $\mathcal{S}(x;B)$ es el conjunto de candidatos generados dentro del presupuesto $B$, $s_{\theta}$ es la puntuación del modelo y $r$ es el verificador/recompensa. Aumentar el compute en tiempo de inferencia expande el conjunto de candidatos o la profundidad de la verificación, pero no repara conocimientos completamente ausentes ni permisos de herramientas faltantes.

\lecturefigure{slide-088.jpg}{Mapa de métodos de post-training e inference-time compute}{CS25 V6 oficial, p.~88}

\teachervoice{00:36:33--00:38:05，el docente considera prompting, fine-tuning, aprendizaje de preferencias, recuperación y uso de agentes/herramientas como diferentes puntos de intervención. Nota de clase: el comportamiento desplegado surge de la interacción entre el modelo base y el sistema de inferencia.}

\begin{importantbox}{Cuatro significados de "mejorar un modelo"}
\begin{enumerate}
  \item \textbf{Actualizar pesos}: SFT, DPO, RL, etc., cambian directamente los parámetros;
  \item \textbf{Modificar decodificación}: autoconsistencia, búsqueda y reranking alteran la selección de candidatos;
  \item \textbf{Añadir contexto/estado}: RAG, memoria y salidas de herramientas añaden información visible;
  \item \textbf{Cambiar el entorno}: el ciclo del agente permite al modelo actuar y obtener nuevas observaciones.
\end{enumerate}
Los costos, la capacidad de reversión y los patrones de fallo de las cuatro son completamente distintos.
\end{importantbox}

\subsection{De CoT a descomposición ejecutable}

Chain-of-Thought (CoT) exige que el modelo produzca pasos intermedios, dividiendo una mapeo único en razonamiento secuencial. Puede mejorar el rendimiento en tareas complejas o generar explicaciones que suenen lógicas pero sean infieles. Tree-of-Thoughts amplía esto a una búsqueda de múltiples ramas; Program-of-Thought delega parte del cálculo a un intérprete; Socratic questioning y graph de cómputo definen explícitamente las dependencias entre subproblemas.

\lecturefigure{slide-089.jpg}{Chain-of-Thought: permitir al modelo desplegar pasos intermedios antes de la respuesta}{CS25 V6 oficial, p.~89}

\lecturefigure{slide-090.jpg}{Comparativa de ejemplos entre prompting estándar y CoT prompting}{CS25 V6 oficial, p.~90}

\lecturefigure{slide-091.jpg}{Tree-of-Thoughts: expansión desde una cadena de razonamiento única a búsqueda ramificada}{CS25 V6 oficial, p.~91}

\lecturefigure{slide-092.jpg}{Program-of-Thought: delegar cómputos ejecutables al intérprete de código}{CS25 V6 oficial, p.~92}

\begin{knowledgebox}{Lectura de la figura: los cuatro métodos alteran interfaces distintas}
CoT añade tokens secuenciales; ToT añade ramas de candidatos y un controlador de búsqueda; Program-of-Thought añade una herramienta ejecutable; Socratic/graph de cómputo añaden una estructura de dependencia explícita. Todos pueden llamarse "razonamiento", pero difieren en compute, verificabilidad y rutas de propagación de errores.
\end{knowledgebox}

Socratic decomposition puede escribirse como un grafo acíclico dirigido $G=(V,E)$, donde cada nodo $v_i$ es un subproblema y la arista $(v_j,v_i)$ indica que $i$ depende de $j$. Solo se resuelve $v_i$ cuando todos sus predecesores están completos. Esta estructura facilita la verificación local, pero introduce error de descomposición: si se omiten nodos clave al inicio, cada paso posterior puede ser localmente correcto y globalmente fallido.

\lecturefigure{slide-093.jpg}{Socratic questioning: preguntas recursivas que dividen problemas complejos en subproblemas}{CS25 V6 oficial, p.~93}

\lecturefigure{slide-094.jpg}{Graphs de cómputo: organizar descomposición y resultados intermedios mediante dependencias}{CS25 V6 oficial, p.~94}

\teachervoice{00:38:05--00:40:36，el docente divide las extensiones de CoT en branching, ejecución y descomposición estructurada. Recordatorio del material: más tokens de razonamiento son presupuesto de búsqueda, no equivalen a que esos tokens reflejen fielmente el proceso causal interno del modelo.}

\begin{lstlisting}[language=Python,caption={Pseudocódigo de búsqueda de razonamiento estructurado con verificador}]
frontier = [initial_state(problem)]
for depth in range(max_depth):
    candidates = []
    for state in frontier:
        for thought in model.propose(state, branches=k):
            next_state = state.extend(thought)
            score = verifier(next_state)
            candidates.append((score, next_state))
    frontier = top_b(candidates, beam_width)
return select_final(frontier)
\end{lstlisting}

\begin{warningbox}{Tres niveles de evidencia para la reasoning trace}
Una trazabilidad puede ser una \textbf{explicación legible}, un \textbf{scratchpad que ayuda a obtener la respuesta} o una \textbf{evidencia causal fiel del cálculo interno}. Los dos primeros son comunes; el tercero requiere intervención, contrafactuales o evidencia mecánica. No se debe afirmar que "el modelo pensó exactamente así" solo porque los pasos textuales sean correctos.
\end{warningbox}

\subsection{RLHF, DPO, RLAIF y GRPO: fuentes de retroalimentación y rutas de optimización distintas}

Reinforcement Learning from Human Feedback (RLHF) típicamente recopila preferencias primero, entrena un modelo de recompensa y luego usa optimización de política para cambiar la estrategia de generación. Direct Preference Optimization (DPO) salta el modelo de recompensa explícito y el RL online, optimizando directamente los pares de preferencias:
\[
\mathcal{L}_{\mathrm{DPO}}
=-\log\sigma\!\left(\beta\left[
\log\frac{\pi_{\theta}(y^+\mid x)}{\pi_{\mathrm{ref}}(y^+\mid x)}
-\log\frac{\pi_{\theta}(y^-\mid x)}{\pi_{\mathrm{ref}}(y^-\mid x)}
\right]\right).
\]
$y^+$ y $y^-$ son respuestas preferidas/rechazadas, $\pi_{\mathrm{ref}}$ es la política de referencia y $\beta$ controla el grado de desviación.

\lecturefigure{slide-096.jpg}{RLHF: preferencias humanas, modelo de recompensa y optimización de política}{CS25 V6 oficial, p.~96}

\lecturefigure{slide-097.jpg}{DPO: actualizar la política directamente desde pares de preferencias sin entrenar explícitamente un modelo de recompensa}{CS25 V6 oficial, p.~97}

RLAIF usa IA para generar o juzgar preferencias, reduciendo costos manuales, pero puede amplificar los sesgos del modelo maestro. GRPO muestrea un conjunto de respuestas bajo el mismo prompt y construye una ventaja relativa dentro del grupo. Si la recompensa es $r_i$, con media y desviación estándar $\mu_r,\sigma_r$ del grupo, se puede escribir como
\[
\hat{A}_i=\frac{r_i-\mu_r}{\sigma_r+\epsilon}.
\]
Evita el costo de un modelo value independiente, pero aún depende de la fiabilidad del verificador/recompensa y puede generar altos costos computacionales mediante muestreo por grupos.

\lecturefigure{slide-098.jpg}{RLAIF: retroalimentación de IA sustituye parcialmente la anotación manual de preferencias}{CS25 V6 oficial, p.~98}

\lecturefigure{slide-099.jpg}{GRPO: recompensa relativa intra-grupo y optimización de política de razonamiento}{CS25 V6 oficial, p.~99}

\begin{knowledgebox}{Términos clave: ¿en qué difieren realmente los métodos de optimización de preferencias?}
\begin{tabularx}{\textwidth}{p{0.17\textwidth}p{0.24\textwidth}X}
\toprule
Método & Fuente de retroalimentación/modelo intermedio & Mecánica central y riesgos principales \\
\midrule
RLHF & Preferencias humanas + modelo de recompensa & Optimización de política online/casi-online; hacking de recompensas y alta complejidad de entrenamiento \\
DPO & Pares de preferencias humanos o IA & Objetivo tipo clasificación directa; más simple pero aún sujeto a la calidad de los datos de preferencia \\
RLAIF & Crítico de IA o retroalimentación constitucional & Reduce costos manuales; el sesgo del maestro puede ser amplificado sistemáticamente \\
GRPO & Recompensa verificable + muestras intra-grupo & Ventaja relativa, sin modelo value independiente; costo de muestreo y sesgo del verificador son evidentes \\
\bottomrule
\end{tabularx}
\end{knowledgebox}

\subsection{Supervisión del proceso: recompensar pasos o solo el resultado}

La supervisión de resultados (Outcome supervision) mira solo la respuesta final; el modelo puede obtener puntos accidentalmente mediante atajos. La supervisión del proceso (process supervision) proporciona etiquetas o recompensas para los pasos intermedios. Si la trazabilidad de razonamiento es $z_{1:m}$, se puede escribir como
\[
R_{\mathrm{process}}(z)=\sum_{i=1}^{m}w_i r_i(z_i,z_{<i},x).
\]
$r_i$ evalúa el paso $i$, y $w_i$ es el peso. La dificultad radica en quién define los "pasos correctos", cómo manejar múltiples rutas equivalentes y si el modelo aprende a complacer al verificador en lugar de resolver el problema.

\lecturefigure{slide-100.jpg}{Supervisión del proceso: pasar de supervisar la respuesta final a supervisar pasos intermedios de razonamiento}{CS25 V6 oficial, p.~100}

\teachervoice{00:40:36--00:43:12，el docente distingue RLHF, DPO, RLAIF, GRPO y supervisión del proceso: difieren en la fuente de retroalimentación, si construyen explícitamente un modelo de recompensa y si recompensan el resultado o los pasos.}

\begin{warningbox}{Las tareas verificables no implican que el verificador sea imparcial}
Las respuestas matemáticas, las pruebas de código y las verificaciones de formato se pueden calificar automáticamente fácilmente, pero la recompensa a menudo cubre solo partes medibles. El modelo puede generar respuestas sobreajustadas a las pruebas, aprovechar tolerancias numéricas o producir "pasos calificables" redundantes. Por tanto, el diseño del verificador es en sí mismo un objetivo de entrenamiento, no una herramienta de medición neutral.
\end{warningbox}

Desde la infraestructura, DPO requiere principalmente pares de preferencias y retropropagación estándar; RLHF/GRPO requieren generación de rollout, cálculo de recompensa, logprob de política/modelo de referencia y pueden intercambiar secuencias de longitud variable masiva entre múltiples máquinas. Los tokens de razonamiento también son costo de entrenamiento: las diferencias en longitud de razonamiento dentro del mismo batch causan padding y carga desigual. Al elegir un algoritmo, se debe comparar simultáneamente la construcción de datos offline, muestreo online, latencia de recompensa, copias múltiples de modelos en memoria y reproducibilidad. Que un algoritmo tenga mayor eficiencia de tokens en un artículo no significa que sea mejor en tiempo real (wall-clock) o uso del clúster.

\subsection{Resumen de la sección}

El post-training puede modificar parámetros, el compute en tiempo de inferencia puede expandir la búsqueda, la recuperación/herramientas pueden añadir información y la supervisión del proceso puede cambiar la granularidad de la recompensa. CoT, ToT, Program-of-Thought, RLHF, DPO, RLAIF y GRPO no son una misma clase de "truco de razonamiento"; operan sobre espacios de candidatos, ejecución externa, datos de preferencia y optimizadores respectivamente.
\section{Agentes que se automejoran: transformar la generación única en un sistema de ciclo cerrado}

Post-entrenamiento, el modelo aún puede verse como una función de entrada a salida; los agentes permiten que el modelo perciba, decida, actúe y actualice su estado repetidamente en un entorno. Lo crucial no es usar el nombre "agente", sino si el sistema posee objetivos, observación, espacio de acción, memoria, permisos de herramientas y condición de terminación.

\subsection{Definición operativa del agente}

Supongamos que el estado del entorno es $s_t$, la observación visible para el modelo es $o_t$, la memoria interna es $m_t$ y la estrategia de acción es
\[
a_t\sim\pi_{\theta}(a\mid o_t,m_t,g),
\]
donde $g$ es el objetivo. El entorno devuelve $o_{t+1}$ y la actualización de memoria es $m_{t+1}=U(m_t,o_{t+1},a_t)$. La capacidad del agente proviene del ciclo cerrado formado por la estrategia, las herramientas, la actualización de estado y la retroalimentación del entorno, no está determinada solo por los parámetros del modelo.

\lecturefigure{slide-102.jpg}{Componentes del agente de IA: objetivo, decisión, herramientas, memoria y retroalimentación}{CS25 V6 official deck, p.~102}

\lecturefigure{slide-103.jpg}{Ciclo del agente: percibir, razonar, actuar, observar, reflexionar y repetir}{CS25 V6 official deck, p.~103}

\teachervoice{00:43:12--00:44:05, el docente comprime la diferencia entre agentes y modelos de una sola vez en un ciclo. Nota de clase: los prompts multironda sin observación o actualización de estado no necesariamente constituyen una interacción ambiental real.}

\begin{importantbox}{Campos mínimos del contrato del agente}
El objetivo $g$, la observación visible, las acciones/herramientas permitidas, la política de memoria, el presupuesto, la condición de parada y los límites de recuperación de errores y permisos deben definirse explícitamente. Si falta cualquiera de estos elementos, la "autonomía" se convertirá en una promesa ambigua no auditable.
\end{importantbox}

\subsection{Automejora: modificar salida, memoria o estrategia}

La "automejora" puede referirse a tres niveles distintos: refinamiento de la respuesta actual; escribir resúmenes de fallos en la memoria episódica; o actualizar parámetros mediante entrenamiento. La diapositiva 104 propone el marco general, mientras que las diapositivas 105--106 muestran el refinamiento y la reflexión por separado. Los dos primeros alteran principalmente el contexto actual o el estado externo, lo cual no equivale a que el modelo haya aprendido permanentemente nuevas capacidades.

\lecturefigure{slide-104.jpg}{Automejora: generación, retroalimentación, revisión y acumulación de experiencia}{CS25 V6 official deck, p.~104}

\lecturefigure{slide-105.jpg}{Refinamiento: crítica e iteración revisada del mismo output}{CS25 V6 official deck, p.~105}

\lecturefigure{slide-106.jpg}{Reflexión: resumir fallos en memoria reutilizable para tareas posteriores}{CS25 V6 official deck, p.~106}

Si la candidata $k$-ésima es $y^{(k)}$, el crítico genera retroalimentación $c^{(k)}$, y el refinamiento puede escribirse como
\[
y^{(k+1)}\sim p_{\theta}(y\mid x,y^{(k)},c^{(k)}).
\]
Solo cuando los errores del crítico y del generador no están completamente correlacionados, la retroalimentación puede localizar problemas reparables y el presupuesto permite múltiples iteraciones, es posible que el ciclo mejore el rendimiento. De lo contrario, el sistema repetirá la reescritura de redacción o amplificará el mismo error.

\subsection{ReAct: alternar razonamiento y acción}

ReAct alterna pensamiento, acción y observación en la trayectoria. La acción llama a búsqueda, bases de datos, código u otras APIs; la observación devuelve los resultados reales al modelo. Tiene más anclaje ambiental que el CoT puro, pero aumenta errores de selección de herramientas, parámetros incorrectos, violación de permisos, inyección de prompts y estados externos no reproducibles.

\lecturefigure{slide-107.jpg}{Ciclo cerrado alternado de razonamiento, acción de herramienta y observación en ReAct}{CS25 V6 official deck, p.~107}

\begin{lstlisting}[language=Python,caption={Ciclo mínimo del agente con permisos y condición de terminación}]
state = initialize(goal, memory)
for step in range(max_steps):
    proposal = model.plan(state)
    if proposal.action == "finish":
        return proposal.answer
    authorize(proposal.action, policy=tool_policy)
    observation = tools.execute(proposal.action)
    state = update_state(state, proposal, observation)
raise BudgetExceeded(state)
\end{lstlisting}

\teachervoice{00:44:05--00:45:43, el docente considera que el refinamiento, la reflexión y ReAct consisten en insertar retroalimentación en diferentes ubicaciones: después del output, memoria entre tareas o después de una acción ambiental. Experiencia práctica: a mayor longitud del ciclo, más se requiere ingeniería para el estado, permisos y terminación.}

\begin{warningbox}{Más iteraciones no garantizan fiabilidad}
Si el crítico y el generador comparten puntos ciegos, la reflexión resumirá errores en memoria más persuasiva; si la observación de las herramientas está contaminada por inyección de prompts, la siguiente razonación continuará desde un estado erróneo. La evaluación del agente debe medir tasa de éxito, costo, capacidad de recuperación y efectos secundarios, no solo la calidad final del texto.
\end{warningbox}

La evaluación del agente debe tomar la trayectoria en lugar de la respuesta final como unidad mínima de auditoría. Se necesita registrar cada llamada a herramienta, parámetros, origen de observación, cambios de estado y causa de parada, luego estadizar éxito de tarea, pasos inválidos, acciones peligrosas, tasa de recuperación, costo token/latencia y violaciones de permisos. Para herramientas irreversibles, también se debe distinguir simulación, ejecución previa a producción y ejecución real. Si el benchmark solo otorga recompensa final, el sistema podría tener éxito mediante atajos accidentales ocultando rutas inseguras; si solo se mira una trazada bonita, podría recompensar explicaciones extensas. El objetivo de un agente confiable es completar tareas de manera estable con costos controlables y detenerse en un estado seguro ante fallos.

\subsection{Resumen de la sección}

El agente transforma el modelo en una estrategia de ciclo cerrado, pero cada interfaz del ciclo introduce nuevas capacidades de observación y patrones de fallo. El refinamiento modifica la candidata actual, la reflexión cambia la experiencia externa y ReAct añade observación real. Para hablar de automejora, es necesario especificar si lo que se modifica es el output, memoria, estrategia de herramientas o parámetros.

\section{Aplicaciones multimodales: unificar el backbone, no los estándares de evidencia}

La interfaz secuencial del Transformer permite que las imágenes en parches (patches), los tokens de texto y las características de redes cerebrales ingresen a backbones similares; sin embargo, la observación, la etiqueta y la evaluación son distintas en cada dominio. El enfoque de esta sección no es enumerar "qué más se puede hacer", sino comparar los adaptadores modales, el objetivo de entrenamiento y los límites de la evidencia científica.

\subsection{Mapa de aplicaciones: un mismo bloque conectando mundos distintos}

El curso presenta asistentes de lenguaje, AlphaFold, modelos médicos, juegos, robots y sistemas multimodales. La expansión de las aplicaciones surge de una combinación en tres capas: codificar las modalidades originales en tokens; aprender representaciones mediante preentrenamiento a gran escala; e incorporar decodificadores (decoders), buscadores (retrievers), políticas (policies) o cabezas específicas del dominio (domain-specific heads) en la capa de tareas. Los bloques Transformer compartidos cubren únicamente la capa de cálculo intermedia.

\lecturefigure{slide-109.jpg}{Mapa de aplicaciones del Transformer: lenguaje, ciencia, medicina, juegos y robots}{CS25 V6 official deck, p.~109}

\begin{knowledgebox}{Lectura de la figura: los logotipos de las aplicaciones no representan el mismo tipo de evidencia}
Las capturas de pantalla demuestran la existencia del sistema; los benchmarks prueban las puntuaciones bajo un protocolo específico; y las conclusiones clínicas o científicas requieren validación externa. Colocar todos los logotipos en una sola página permite mostrar el alcance del impacto, pero no se puede generalizar su riesgo y eficacia con un único concepto de "precisión" (accuracy).
\end{knowledgebox}

\subsection{Vision Transformer: reescribiendo la imagen como una secuencia de parches}

Dada una imagen de dimensiones $H\times W\times C$ y un tamaño de parche ($P$), el número de parches es:
\[
N=\frac{H}{P}\frac{W}{P}.
\]
Cada parche se aplanó (flatten) y proyectó en dimensión $d$, más lo cual se le suma la representación posicional antes de enviarlo al Transformer. Un $P$ menor ofrece una granularidad más fina, pero el aumento de $N$ eleva el costo del mecanismo de atención proporcionalmente a $N^2$. Las CNN inyectan fuertes sesgos previos (prior knowledge) mediante localidad y compartición de pesos; en cambio, los ViT tienen un sesgo previo más débil y suelen depender más de la escala de datos y del aumento de datos (augmentation).

\lecturefigure{slide-111.jpg}{ViT frente a CNN: contraste entre el sesgo previo de convolución local y la interacción global de tokens}{CS25 V6 official deck, p.~111}

\subsection{CLIP: alineando imagen y texto mediante un objetivo contrastivo}

CLIP codifica por separado la imagen y el texto, haciendo que los pares coincidentes sean similares y los no coincidentes se separen dentro de un batch. Para una representación de imagen $u_i$ y una de texto $v_j$, la similitud puede expresarse como $s_{ij}=u_i^{\top}v_j/\tau$, aplicando después una entropía cruzada (cross-entropy) por filas y columnas. Aprende un espacio conjunto (joint space) transferible, pero también incorporan sesgos de descripción, propiedad intelectual y prejuicios sociales presentes en los datos de entrenamiento.

\lecturefigure{slide-112.jpg}{CLIP: alineación del codificador de imagen y el de texto en un espacio compartido}{CS25 V6 official deck, p.~112}

\teachervoice{00:45:43--00:48:30, el docente utiliza ViT y CLIP para ilustrar la premisa del backbone unificado: primero convertir las modalidades en tokens o embeddings, y luego seleccionar el objetivo adecuado. También enfatiza que el sesgo inductivo débil de los ViT suele requerir conjuntos de datos más grandes.}

\begin{importantbox}{Unificar el backbone no es lo mismo que unificar el modelo}
La visión requiere parches/tokenizadores, el texto necesita un tokenizador de subpalabras (subword), y la fMRI necesita representaciones de regiones cerebrales y series temporales. El codificador, el aumento de datos, la función de pérdida (loss) y la cabeza (head) pueden diferir. Lo que se denomina "Transformer multimodal" debe especificar qué parámetros comparten, en qué capa ocurre la fusión y cómo se alinean las distintas modalidades.
\end{importantbox}

\subsection{Modelo fundacional de fMRI: construyendo relaciones de redes cerebrales como un problema de representación}

La resonancia magnética funcional (fMRI) mide indirectamente la actividad cerebral mediante señales de dependencia del nivel de oxígeno en sangre. Divisiones como las de Yeo networks dividen el cerebro en redes funcionales como visión, atención y modo por defecto. Los modelos pueden codificar las series temporales de regiones o redes en tokens para luego aprender interacciones entre redes. El objetivo aquí es el aprendizaje de representaciones y la comparación de cohortes, no la lectura directa del "pensamiento".

\lecturefigure{slide-114.jpg}{fMRI y Yeo networks: desde señales volumétricas hasta representaciones de redes funcionales}{CS25 V6 official deck, p.~114}

\lecturefigure{slide-115.jpg}{Representaciones de fMRI fundamentadas (Grounded): tokens de redes, atención e interfaces interpretables}{CS25 V6 official deck, p.~115}

Si los tokens de la red son $X\in\mathbb{R}^{n\times d}$, la atención entre redes puede expresarse aún como $\operatorname{softmax}(QK^\top/\sqrt{d_k})V$, pero el significado semántico de cada token es una red funcional y no una palabra. El valor de las representaciones fundamentadas radica en asociar la atención con regiones cerebrales conocidas; sin embargo, su limitación es que los pesos de atención no son equivalentes a la conectividad causal.

\subsection{Dependencia entre redes y evidencia longitudinal}

La diapositiva 117 compara gráficos de dependencia de red entre cohortes control, MCI y AD, mientras que la 118 muestra cambios longitudinales y comparaciones de grupos. Al interpretar estos gráficos, primero se deben revisar la definición de la cohorte, el tamaño de la muestra y los intervalos de confianza antes de observar los enlaces de la red; las diferencias en enfermedades pueden verse afectadas por la edad, el protocolo de escaneo y artefactos de movimiento. La correlación del modelo puede generar hipótesis biológicas, pero las conclusiones clínicas requieren cohortes independientes y validación mecanística.

\lecturefigure{slide-117.jpg}{Dependencia entre redes: contraste de relaciones en redes para control, MCI y AD}{CS25 V6 official deck, p.~117}

\lecturefigure{slide-118.jpg}{Cartografía longitudinal: representaciones de redes y comparaciones entre grupos según edad o progresión de la enfermedad}{CS25 V6 official deck, p.~118}

\teachervoice{00:48:30--00:52:47, Karán explica que la fMRI es una señal sustituta del oxígeno y limita los resultados a representaciones de redes y asociaciones a nivel de cohorte. Nota de clase: las arcos de atención o dependencia en los gráficos no indican dirección causal neuronal.}

\begin{warningbox}{Tres niveles de incertidumbre en modelos de neuroimagen}
Capa de medición: la señal BOLD es un sustituto indirecto de la actividad neuronal; capa de datos: la muestra, el movimiento y el protocolo de escaneo pueden introducir sesgos; capa del modelo: las diferencias en atención o embeddings pueden tener múltiples interpretaciones. Solo la replicación a través de cohortes, el control de variables clínicas y los experimentos externos pueden llevar la evidencia de representación hacia mecanismos biológicos.
\end{warningbox}

Un flujo de trabajo científico más completo consiste en: aprender representaciones primero en una cohorte de entrenamiento, realizar predicciones en sujetos reservados (held-out subject); luego realizar un análisis de sensibilidad con covariables como sitio, edad, sexo y movimiento para verificar si las diferencias entre redes son estables; replicar con una cohorte externa; y finalmente discutir las implicaciones clínicas. Si el modelo produce una trayectoria longitudinal, también debe reportarse la incertidumbre y cómo se manejan las visitas faltantes. El Transformer puede integrar series temporales de alta dimensión, pero no corrige automáticamente los sesgos de selección en estudios observacionales. El valor de incluir estos casos en el resumen del curso es precisamente advertir al lector que "multimodal" implica no solo cambios en el formato de entrada, sino también una actualización de los estándares de evidencia según el dominio.

\subsection{Resumen de la sección}

El valor unificador del Transformer reside en el cálculo secuencial combinable, no en eliminar las diferencias entre modalidades. ViT utiliza tokens de parches, CLIP usa alineación contrastiva y los modelos de fMRI emplean tokens de redes funcionales. Cada aplicación debe redefinir la observación, el objetivo, la métrica y la intensidad de la evidencia; un backbone compartido no sustituye la validación del dominio.


\section{Después del escalado: capacidades ausentes, eficiencia y alucinación}

El crecimiento de las aplicaciones no ha eliminado los problemas fundamentales. La clase divide el "futuro" en dos categorías: capacidades que pueden seguir siendo ingenierizadas, como la eficiencia on-device, la personalización y el monitoreo; y problemas cuya definición aún es poco clara, como la comprensión del mundo, las alucinaciones y la adaptación a largo plazo. Al leer este capítulo, distinga entre las tendencias empíricas de las leyes de escalado, los deseos de producto y los mecanismos verificables.

\subsection{Próximos pasos y brechas: un modelo más grande es solo una dirección}

La diapositiva 121 enumera el crecimiento visible en aplicaciones: combinaciones de modelos base (foundation models), agentes, robótica, personalización, salud y entretenimiento interactivo. La diapositiva 122 lista los ingredientes faltantes: razonamiento confiable, comprensión causal, aprendizaje continuo, alineación y eficiencia. Los primeros responden a "¿dónde más se pueden desplegar?", mientras que los segundos responden a "¿qué falta para un despliegue confiable?".

\lecturefigure{slide-121.jpg}{Aplicaciones futuras: agentes, robótica, personalización, salud y sistemas interactivos}{CS25 V6 oficial deck, p.~121}

\lecturefigure{slide-122.jpg}{Brechas futuras: razonamiento, causalidad, multimodalidad, aprendizaje continuo y alineación}{CS25 V6 official deck, p.~122}

\teachervoice{00:52:47--00:53:34，el docente primero enumera aplicaciones potenciales y luego se dirige rápidamente a los ingredientes faltantes. Nota de clase: la lista de aplicaciones define una dirección; la lista de brechas define los problemas de investigación y los riesgos de despliegue.}

\begin{knowledgebox}{Convertir el slide futuro en un problema de ingeniería}
La "personalización" requiere políticas de identidad/memoria; la "robótica" requiere un bucle percepción-acción y restricciones de seguridad; la "salud" requiere calibración, procedencia y validación clínica; el "aprendizaje continuo" requiere estados actualizables y control del olvido. Solo al especificar entradas, actualizaciones y métricas, las visiones futuras pasan de ser eslóganes a convertirse en planes de investigación.
\end{knowledgebox}

\subsection{Minificado/on-device: los modelos deben entrar en hardware limitado}

Los LLM on-device están sujetos a restricciones de memoria, consumo energético, latencia y canales de actualización. La compresión puede dividirse en cuantización, podado (pruning), destilación, adaptación de rango bajo (low-rank adaptation) y rediseño arquitectónico. Si el número de parámetros es $N$ y la anchura de bits por parámetro es $b$, el almacenamiento de pesos es aproximadamente
\[
M_{\mathrm{weight}}\approx \frac{Nb}{8}\ \text{bytes}.
\]
El pico real también incluye activaciones, caché KV, búferes de tiempo de ejecución y tokenizador/embedding. Cuantizar solo los pesos a 4-bits no significa que todo el sistema se reduzca cuatro veces.

\lecturefigure{slide-123.jpg}{LLMs minificados y on-device: objetivos de eficiencia, privacidad y usabilidad}{CS25 V6 oficial deck, p.~123}

\begin{importantbox}{Frontera de Pareto de optimización en el lado del cliente}
Se deben comparar simultáneamente calidad, latencia del primer token, latencia por token, memoria pico, consumo energético, diseño térmico y actualizabilidad. Un modelo con puntuaciones similares en un benchmark pero que se calienta continuamente o no puede cargar contextos largos no constituye una solución viable para el lado del cliente.
\end{importantbox}

\subsection{Beneficios decrecientes del escalado y los ejes del siguiente lote}

La ley empírica de escalado suele escribirse como
\[
L(N,D,C)\approx L_{\infty}+aN^{-\alpha}+bD^{-\beta}+cC^{-\gamma},
\]
donde $N$ son los parámetros, $D$ los datos y $C$ la capacidad de cómputo; los exponentes describen las mejoras marginales en un régimen específico. Esta relación predice la pérdida promedio, pero no garantiza que mejoren simultáneamente la confiabilidad fáctica, el uso de herramientas o el razonamiento a largo plazo (long-horizon reasoning). A medida que los modelos crecen, aumentan los costos de entrenamiento y servicio, y se vuelve más difícil controlar la calidad de los datos y la contaminación en las evaluaciones.

\lecturefigure{slide-124.jpg}{Límites del escalado: beneficios decrecientes, cuellos de botella de datos y costos de cómputo}{CS25 V6 oficial deck, p.~124}

\lecturefigure{slide-125.jpg}{Después del escalado: nuevas arquitecturas, procedimientos de aprendizaje y direcciones teóricas}{CS25 V6 official deck, p.~125}

\teachervoice{00:53:34--00:55:23，el docente señala que el simple escalado ya presenta rendimientos decrecientes (diminishing returns) y enumera currículo, procedimiento de razonamiento, explicación teórica y arquitecturas alternativas como los ejes del siguiente lote. La guía indica: esto no significa que "el escalado haya dejado de funcionar", sino que los costos marginales han aumentado.}

\begin{warningbox}{La pérdida (loss) de escalado y el escalado de capacidad son diferentes}
Una menor entropía cruzada puede mejorar la generación promedio, pero no garantiza una mejora monótona en habilidades raras, confiabilidad fáctica o seguridad del agente. La capacidad a menudo depende de la cobertura de datos, elicitation, entrenamiento posterior (post-training) y umbrales de evaluación; no se puede deducir un cronograma para todas las capacidades desde una sola curva de pérdida.
\end{warningbox}

\subsection{Bases de la alucinación: no pongamos todos los errores en el mismo cubo por ahora}

La clase primero enumera significados comunes: generación de contenido no real, conflicto con las fuentes, falta de soporte, desinformación o acciones erróneas. El problema radica en que diferentes tareas utilizan diferentes referencias: los resúmenes tienen un documento fuente, las preguntas abiertas tienen hechos reales y los agentes también tienen observaciones y consecuencias de la acción. Si la definición no especifica la referencia y la información visible para el modelo, no se pueden comparar los benchmarks ni distinguir entre falta de conocimiento, falta de observación o fallo en la planificación.

\lecturefigure{slide-127.jpg}{Bases de la alucinación: información errónea, falta de soporte y salida engañosa}{CS25 V6 official deck, p.~127}

\lecturefigure{slide-128.jpg}{Por qué es importante la alucinación: confiabilidad, dominio profesional y confianza}{CS25 V6 official deck, p.~128}

\teachervoice{00:55:23--00:58:11，el docente usa resumido, preguntas (QA) y agentes para explicar cómo la misma salida puede tener diferentes etiquetas en diferentes configuraciones. Nota de clase: defina primero la referencia y luego discuta las mitigaciones.}

\subsection{Sin una definición unificada: perspectiva del modelo del mundo}

El marco unificado citado en la clase explica la alucinación como un error de modelado del mundo, no como todos los comportamientos erróneos. Sea $W$ el mundo de referencia y $V$ el mundo visible para el modelo; sea $P$ la política de manejo de conflictos. El conjunto de proposiciones que produce el modelo es $B_{\theta}(V,P)$, y la condición de alucinación se puede escribir como
\[
\exists b\in B_{\theta}(V,P):\quad b\not\models W,
\]
es decir, el modelo forma o expresa una creencia incompatible con el mundo de referencia. Si la creencia es correcta pero el plan de acción es malo, entonces se trata de un error de planificación y no debe llamarse automáticamente alucinación.

\lecturefigure{slide-130.jpg}{Sin definición unificada: los rangos de cobertura de las definiciones actuales de alucinación difieren}{CS25 V6 official deck, p.~130}

\lecturefigure{slide-131.jpg}{Insight central: la alucinación es un error del modelo del mundo, no todos los errores de salida}{CS25 V6 official deck, p.~131}

\lecturefigure{slide-132.jpg}{Marco unificado: mundo de referencia $W$, mundo visible $V$ y política $P$}{CS25 V6 official deck, p.~132}

\teachervoice{00:58:11--01:00:22，el docente comprime la dependencia de la definición en $W, V, P$: ¿qué es real, qué vio el modelo y cómo se resuelven los conflictos. La guía indica: esta formalización expone las suposiciones, pero no resolverá automáticamente las disputas sobre la fuente de la verdad (truth source).}

\begin{knowledgebox}{Uso por primera vez: ¿a qué se refiere world model aquí?}
El modelo del mundo aquí no se refiere exclusivamente a generadores de video, sino al sistema interno de creencias/representaciones sobre los estados, relaciones y leyes relevantes del mundo. La alucinación se limita al conflicto entre esa representación y el mundo de referencia. Si el modelo nunca observó la información, el problema podría ser la observabilidad; si la observación es correcta pero falla la estrategia, el problema podría ser la planificación.
\end{knowledgebox}

\subsection{Unificación a través de configuraciones, pero sin homogeneizar los tipos de error}

La diapositiva 133 mapea configuraciones como resumido, QA y agentes al mismo marco; la diapositiva 134 enfatiza que no todos los errores son alucinaciones. El valor de la unificación radica en clarificar la referencia, la evidencia visible y la regla de conflicto, no en forzar la fusión de los benchmarks. Las funciones de pérdida, observabilidad y estándares de tolerancia a fallos siguen siendo diferentes para cada configuración.

\lecturefigure{slide-133.jpg}{Unificación de configuraciones: mapeo $W,V,P$ para resumen, QA y agentes}{CS25 V6 official deck, p.~133}

\lecturefigure{slide-134.jpg}{No todos los errores son alucinaciones: distinción entre planificación, recompensa y error de creencia}{CS25 V6 official deck, p.~134}

\begin{warningbox}{Una misma etiqueta de error puede ocultar tres cadenas de causa raíz}
\textbf{Error de conocimiento}: las reglas correctas no están en los parámetros; \textbf{Error de observación}: la evidencia correcta no entra en el contexto o falla la recuperación; \textbf{Error de decisión}: la creencia es correcta pero la planificación/acción es errónea. Los tres requieren mitigaciones distintas: datos de entrenamiento, interfaces de recuperación y planners/verificadores no pueden sustituirse entre sí.
\end{warningbox}

\subsection{Por qué importa la definición: evaluación, intervención y entornos sintéticos}

Una vez definidos $W, V, P$, se puede construir un entorno controlado: fijar los valores verdaderos, cambiar la evidencia visible y especificar reglas de conflicto para generar sistemáticamente casos de alucinación. Así se pueden comparar RAG, calibración, abstención, auto-revisión (self-check) y verificación de herramientas. Sin embargo, los benchmarks sintéticos solo cubren el entorno modelado y aún pueden omitir el mundo abierto, las permisiones y la actualidad en la realidad.

\lecturefigure{slide-135.jpg}{Función de la definición I: exponer suposiciones y comparar diferentes métodos}{CS25 V6 official deck, p.~135}

\lecturefigure{slide-136.jpg}{Función de la definición II: construir entornos controlados y casos sintéticos a gran escala}{CS25 V6 official deck, p.~136}

\lecturefigure{slide-137.jpg}{Visión general de la alucinación: modelado del mundo, configuración y evaluación}{CS25 V6 official deck, p.~137}

\teachervoice{01:00:22--01:03:52，el docente enfatiza el valor de ingeniería de una definición formal: permite comparar configuraciones, generar casos sintéticos, entrenar y evaluar mitigaciones; además, todas las conclusiones dependen de cómo se modelan los valores verdaderos y el mundo visible.}

\begin{importantbox}{Orden de diagnóstico para la mitigación de alucinaciones}
Primero confirme si el mundo de referencia es confiable; luego verifique si la evidencia entró en $V$; después juzgue si la creencia del modelo es errónea; finalmente revise el planificador/acción. Si al principio se le pide al modelo que "piense de nuevo", solo podría estar generando explicaciones más largas sobre evidencia errónea.
\end{importantbox}

Por ejemplo, si un agente médico recomienda un medicamento incorrecto, hay al menos cuatro causas raíz diferentes: la base de conocimientos en sí misma está obsoleta (referencia $W$ errónea); la recuperación no trajo el historial de alergias ($V$ incompleta); el modelo interpretó mal la evidencia visible (creencia del modelo del mundo errónea); la creencia es correcta pero la política de acción ignora las contraindicaciones (error de planificación/alineación). Las reparaciones para estas cuatro situaciones son: actualizar la fuente, mejorar recuperación/permisos, entrenar o calibrar el modelo, y modificar el planificador con reglas de seguridad. Etiquetar todas como alucinación haría que el análisis de causa raíz fallara y haría que los benchmarks parezcan unificados pero sean inaplicables en la práctica.

\subsection{Resumen de la sección}

El escalado sigue siendo efectivo, pero no resuelve automáticamente la eficiencia, la adaptación a largo plazo ni la confiabilidad. El despliegue en el lado del cliente requiere una contabilidad completa de recursos; las leyes de escalado solo describen regímenes específicos de pérdida. Para la alucinación, primero hay que definir claramente $W, V, P$ y separar los errores de creencia de los errores de observación/planificación; una definición precisa es la entrada para el diagnóstico y la evaluación, no la respuesta en sí misma.


\section{Memory, continual learning, interpretability y alignment}

El capítulo anterior discutió las representaciones erróneas de los modelos sobre el mundo actual. Este capítulo indaga cómo los sistemas guardan información a lo largo del tiempo, actualizan sus parámetros y mantienen la coherencia de objetivos en nuevos entornos. Memory, continual learning, interpretability y alignment suelen aparecer juntos porque abordan conjuntamente "en qué se convierte un sistema tras operar por periodos prolongados".

\subsection{Long-term memory: no es solo extender el contexto}

La memoria a largo plazo incluye al menos escritura, recuperación, actualización, resolución de conflictos, compresión, olvido y permisos. Puede denominar a store de memoria $M_t$, nueva observación como $o_t$ y estrategia de actualización como
\[
M_{t+1}=U(M_t,o_t;\rho,\kappa),
\]
donde $\rho$ es la estrategia de retención/eliminación y $\kappa$ las reglas de resolución de conflictos. Si un hecho nuevo contradice uno antiguo, el sistema puede sobrescribirlo, conservar versiones, reducir la confianza o solicitar confirmación; guardar "todo" sin una política genera ruido y riesgos de privacidad.

\lecturefigure{slide-139.jpg}{Objetivo de la memoria a largo plazo: almacenamiento, recuperación y personalización entre sesiones}{CS25 V6 oficial deck, p.~139}

\lecturefigure{slide-140.jpg}{Métodos de memoria: bases vectoriales, resúmenes, compresión, actualización y olvido}{CS25 V6 oficial deck, p.~140}

\teachervoice{01:03:52--01:06:32, el docente enfatiza que la dificultad de memory no radica en "si se puede guardar", sino en contradicción, pruning y relevancia. Experiencia práctica: un contexto sin límites no equivale a un gestor de memoria confiable.}

\begin{lstlisting}[language=Python,caption={Pseudocódigo de actualización de memoria con estrategia de conflicto y eliminación}]
def update_memory(memory, observation, policy):
    candidates = retrieve_conflicts(memory, observation)
    resolution = policy.resolve(observation, candidates)
    memory.apply(resolution)
    if memory.size > policy.capacity:
        memory.evict(policy.rank_for_eviction(memory))
    return memory
\end{lstlisting}

\begin{knowledgebox}{Digestión de términos de Memory}
\textbf{Context memory} son los tokens presentes en el prompt actual; \textbf{episodic memory} guarda eventos pasados; \textbf{semantic memory} almacena hechos o conceptos; \textbf{vector store} es la implementación de recuperación; \textbf{parameter memory} representa las regularidades estadísticas escritas en los pesos. Difieren en velocidad de actualización, trazabilidad del origen y mecanismos de olvido.
\end{knowledgebox}

\subsection{Aprendizaje a lo largo de la vida: ¿cuándo se considera que un modelo realmente aprendió?}

El curso adopta una definición estricta: RAG, context memory y distillation pueden mejorar el uso continuo, pero si los parámetros del modelo y su estructura de habilidades no se adaptan en línea, no necesariamente se llama true continual learning. Los humanos pueden actualizarse mediante interacción continua, mientras que los LLM suelen depender de reentrenamiento offline periódico. Las métricas clave incluyen plasticity, retention, forward transfer, backward transfer y catastrophic forgetting.

\lecturefigure{slide-141.jpg}{Aprendizaje a lo largo de la vida como los humanos: objetivos y dificultades tras el despliegue}{CS25 V6 oficial deck, p.~141}

\lecturefigure{slide-142.jpg}{Trabajo actual vs. necesario: prompt memory, distillation, edición de modelos y adaptación continua real}{CS25 V6 oficial deck, p.~142}

\teachervoice{01:06:32--01:08:08, el docente distingue claramente entre memory/RAG y continual learning a nivel de parámetros. Nota de clase: colocar información en una base vectorial puede actualizar la fuente de conocimiento, pero no altera la dinámica de aprendizaje del modelo en sí.}

\subsection{Edición de modelos: ¿por qué los cambios locales de pesos causan efectos colaterales?}

Métodos como Rank-One Model Editing (ROME) intentan modificar hechos específicos mediante actualizaciones de rango bajo. Abstractamente, si los pesos de una capa son $W$, la actualización es
\[
W'=W+uv^{\top},
\]
donde $u,v$ controlan la dirección de escritura. Un rango bajo local no equivale a semántica local: un mismo parámetro participa en múltiples contextos, por lo que modificar una relación puede afectar paráfrasis, hechos conexos o comportamientos irrelevantes. La evaluación requiere eficacia, generalización, especificidad y portabilidad, no solo observar el prompt objetivo.

\lecturefigure{slide-143.jpg}{Edición de modelos para aprendizaje continuo: escritura de hechos local y sus efectos secundarios}{CS25 V6 oficial deck, p.~143}

\begin{warningbox}{La edición de modelos no es una UPDATE de base de datos}
Los pesos son representaciones distribuidas sin "ranuras" de hechos independientes. Aunque se corrija la pregunta objetivo, también podría destruir conceptos adyacentes, ser válida solo para una formulación específica o entrar en conflicto con el entrenamiento posterior. Los sistemas de edición deben registrar provenance, versiones y rollback, no considerar un éxito único como aprendizaje seguro en línea.
\end{warningbox}

Los sistemas a largo plazo también deben definir la prioridad entre memory y actualización de parámetros. La información de alta vigencia, revocable y con trazabilidad es más adecuada para store externos; las regularidades de uso repetido, estables y ampliamente transferibles son las que merecen escribirse en los parámetros. Antes de actualizaciones en línea, se debe validar locality, retention y seguridad en un shadow model, y luego publicar mediante rollout versionado. Si retrieval, memory y pesos guardan versiones distintas de hechos, el sistema debe establecer quiénes tienen autoridad ante conflictos y exponer el origen al usuario. Por ello, la dificultad del aprendizaje continuo no es un solo paso de gradientes, sino el ciclo completo: recopilación, validación, escritura, monitoreo, rollback y olvido.

\subsection{Interpretability: entender los mecanismos internos más allá de las salidas}

Interpretability intenta identificar qué representa el modelo, qué circuitos generan la salida y si los fallos pueden detectarse anticipadamente. Behavioral eval solo observa entradas y salidas; mechanistic interpretability analiza activación, características, atención/circuitos e intervención. La evidencia fuerte proviene del cambio en las salidas tras modificar variables internas según predicción, no solo de visualizaciones correlacionales.

\lecturefigure{slide-145.jpg}{Interpretabilidad de LLM: desde caja negra hasta feature, circuitos y detección de fallos}{CS25 V6 oficial deck, p.~145}

\begin{importantbox}{Gradiente de evidencia explicativa}
Visualizar atención es descripción; probe es evidencia de decodificabilidad; activación patching/ablation es evidencia de intervención; la replicación entre muestras y modelos respalda afirmaciones mecanísticas más fuertes. Los métodos explicativos también pueden omitir cómputo distribuido, redundante o dependiente del contexto.
\end{importantbox}

\subsection{Problema de alignment: brechas entre objetivos, comportamiento y nuevos entornos}

Alignment exige que los sistemas persigan objetivos correctos y cumplan restricciones en distintos entornos. Formalmente, la recompensa de entrenamiento $R_{\mathrm{train}}$ es solo un proxy para la intención humana $U$. Si una política maximiza
\[
\pi^{\star}=\arg\max_{\pi}\mathbb{E}[R_{\mathrm{train}}],
\]
pero $R_{\mathrm{train}}\neq U$, pueden surgir reward hacking, shortcut, deception o distribution shift. El problema no es simplemente que el modelo "no obedece", sino que la representación del objetivo y la cobertura de supervisión son incompletas.

\lecturefigure{slide-146.jpg}{El problema de alignment: diferencia entre objetivos de optimización del modelo y valores humanos}{CS25 V6 oficial deck, p.~146}

\lecturefigure{slide-147.jpg}{Modos de fallo: reward hacking, deception, shortcut y distribution shift}{CS25 V6 official deck, p.~147}

\teachervoice{01:08:58--01:11:00, el docente divide el riesgo de alignment en shortcut, reward hacking, deception, post-hoc explanation y desajuste en nuevos entornos. Nota de clase: "comportamiento correcto" en el conjunto de entrenamiento no implica que el objetivo se haya interiorizado correctamente.}

\subsection{Técnicas de alignment: diferentes métodos restringen distintos canales de fallo}

Human/AI feedback ajusta preferencias de salida; process supervision verifica pasos intermedios; constitutional AI genera críticas y revisiones mediante reglas; scalable oversight intenta que supervisión limitada cubra sistemas más potentes; interpretability busca mecanismos internos peligrosos. Son complementarias, no sustitutos: las reglas pueden entrar en conflicto, el feedback de IA puede heredar sesgos, las etiquetas de proceso pueden recompensar formato superficial e interpretabilidad también puede ser incompleta.

\lecturefigure{slide-148.jpg}{Técnicas de alignment I: feedback, preferencia, interpretabilidad y process supervision}{CS25 V6 oficial deck, p.~148}

\lecturefigure{slide-149.jpg}{Técnicas de alignment II: constitutional AI, scalable oversight y human feedback}{CS25 V6 oficial deck, p.~149}

\teachervoice{01:11:00--01:13:02, el docente coloca constitutional AI, scalable oversight, process supervision e interpretabilidad en la misma caja de herramientas, indicando explícitamente que el human feedback por sí mismo no garantiza alignment.}

\begin{knowledgebox}{Matriz de métodos de alignment}
\begin{tabularx}{\textwidth}{p{0.2\textwidth}p{0.25\textwidth}X}
\toprule
Método & Objeto principal de restricción & Problemas sin resolver \\
\midrule
Preference feedback & Clasificación y estilo de salida & Representatividad de preferencias, reward hacking \\
Process supervision & Pasos intermedios & Múltiples rutas equivalentes, costo de etiquetado, adaptación superficial \\
Constitutional AI & Principios explícitos & Conflictos de reglas, brecha entre explicación y ejecución \\
Scalable oversight & Supervisión de sistemas potentes & Cómo un supervisor débil juzga fiablemente tareas complejas \\
Interpretability & Feature/circuit internos & Cobertura incompleta, explicaciones engañosas \\
\bottomrule
\end{tabularx}
\end{knowledgebox}

La evaluación de alignment también requiere estratificación. Capability eval mide si el sistema puede completar tareas; alignment eval verifica si sigue eligiendo comportamientos permitidos ante conflictos de objetivos, inducción, límites de permisos y nuevas distribuciones; no son sustituibles entre sí. El conjunto de pruebas debe incluir solicitudes normales, adversarial prompt, herramientas que retornan contenido malicioso, tentaciones locales en trayectorias largas y escenarios de shutdown/abstention. También se debe reportar false refusal, ya que rechazar todas las solicitudes difíciles puede mejorar ciertas métricas de seguridad pero dañar la utilidad. Finalmente, es necesario trazar el frente de Pareto helpfulness--safety--cost, no ocultar compromisos con una única "puntuación de seguridad".

\subsection{Resumen de la sección}

La gestión de memory controla estados externos; continual learning modifica adaptación a largo plazo; model editing intenta escritura local de pesos; interpretabilidad provee evidencia interna; alignment restringe objetivos y comportamientos. No son sustituibles: un sistema con base de recuperación no necesariamente aprende continuamente, y uno que explica la atención no es necesariamente seguro.

\section{Más allá de los Transformers: JEPA, modelos mundiales y modelos de espacio de estados}

En la última parte del curso no se anunció el fin de la era de los Transformers, sino que se propusieron dos direcciones alternativas o complementarias: los modelos mundiales (world models) y JEPA cambian el objeto de predicción, pasando de tokens a estados latentes abstractos; los modelos de espacio de estados (SSM) cambian el cálculo de secuencias, pasando de atención global explícita a actualizaciones de estado recursivas. Ambos se dirigen a objetivos distintos de representación y complejidad computacional.

\subsection{¿Por qué trascender la predicción del siguiente token?}

Los Transformers dominan en múltiples campos, pero aún enfrentan limitaciones en eficiencia, ejecución de contextos largos y comprensión estructural del mundo. En la diapositiva 151, las direcciones emergentes se dividen en modelos mundiales y SSMs. Aquí el término "más allá" (beyond) debe leerse como expansión del espacio de diseño, no como una negación de la atención: muchos sistemas prácticos combinan Transformers, SSM, recuperación (retrieval) y simuladores.

\lecturefigure{slide-151.jpg}{Más allá de los Transformers: brechas en eficiencia, contexto largo y comprensión del mundo}{CS25 V6 oficial deck, p.~151}

\teachervoice{01:12:51--01:13:20，el docente señala que los Transformers aún dominan, pero tienen limitaciones en eficiencia, razonamiento de contexto largo (long-context reasoning) y comprensión del mundo. Nota de clase: los sistemas de próxima generación probablemente serán combinaciones, no una arquitectura única que reemplace por completo a las anteriores.}

\subsection{Modelos mundiales y JEPA: predecir latentes, no reconstruir todos los detalles}

La Arquitectura Predictiva de Representación Conjunta (Joint Embedding Predictive Architecture, JEPA) predice representaciones objetivo en el espacio latente. Sea $f_{\theta}$ el codificador de contexto, $f_{\xi}$ el codificador objetivo y $g_{\phi}$ el predictor; entonces se puede optimizar
\[
\mathcal{L}_{\mathrm{JEPA}}
=\left\|g_{\phi}(f_{\theta}(x_{\mathrm{context}}))-\operatorname{sg}\big(f_{\xi}(x_{\mathrm{target}})\big)\right\|_2^2.
\]
$\operatorname{sg}$ indica stop-gradient (gradiente detenido). El objetivo es predecir representaciones abstractas relacionadas con regiones futuras o faltantes, en lugar de reconstruir pixel a pixel todos los detalles impredecibles. La dificultad radica en evitar el colapso de la representación y asegurar que el latente conserve los estados necesarios para la planificación.

\lecturefigure{slide-152.jpg}{Modelos mundiales y JEPA: cambio desde la predicción del siguiente token hacia representaciones estructurales del entorno}{CS25 V6 oficial deck, p.~152}

\lecturefigure{slide-153.jpg}{Mecanismo JEPA: codificador de contexto, predicción latente y objetivo abstracto}{CS25 V6 oficial deck, p.~153}

\teachervoice{01:13:02--01:14:38，el docente describe JEPA como una dirección para aprender la estructura del entorno y los estados futuros, en lugar de "predecir más píxeles". Recordatorio del material de clase: si la predicción latente forma un modelo mundial útil para la planificación aún requiere evidencia empírica de tareas específicas.}

\begin{warningbox}{Un modelo mundial no se valida solo con su nombre}
El hecho de que un modelo pueda predecir latentes de video, el siguiente estado de observación o token no implica que posea un modelo del mundo causal, intervenible y planificable. Al menos debe verificarse la suficiencia del estado (state sufficiency), la consistencia contrafactual (counterfactual consistency), los despliegues a largo plazo (long-horizon rollout) y las predicciones condicionadas a acciones (action-conditioned prediction).
\end{warningbox}

\subsection{SSM: reemplazar interacciones cuadráticas explícitas con actualizaciones de estado}

Un Modelo de Espacio de Estados (State Space Model, SSM) escribe la secuencia como actualizaciones de estado:
\[
h_t=\bar{A}_t h_{t-1}+\bar{B}_t x_t,
\qquad y_t=C_t h_t+D x_t.
\]
$h_t$ es el estado comprimido; $\bar{A}_t,\bar{B}_t,C_t$ pueden ser fijos o seleccionados según la entrada. Los SSM selectivos (como Mamba) hacen que parte de los parámetros dependa de la entrada para decidir qué escribir, retener y leer. El entrenamiento puede aprovechar técnicas de barrido/convolución en paralelo, mientras que la inferencia actualiza el estado típicamente con complejidad lineal respecto a la longitud de la secuencia.

\lecturefigure{slide-154.jpg}{Modelo de Espacio de Estados: estados continuos, tiempo lineal y ejemplo Mamba}{CS25 V6 oficial deck, p.~154}

\lecturefigure{slide-155.jpg}{Compensaciones SSM: eficiencia, secuencias largas frente a flexibilidad de atención}{CS25 V6 oficial deck, p.~155}

\teachervoice{01:14:38--01:16:20，el docente enfatiza que los SSM son atractivos por su longitud de secuencia y complejidad lineal, pero la atención sigue siendo superior en tareas que requieren direccionamiento flexible del contenido; las conclusiones actuales no están cerradas.}

\begin{knowledgebox}{Términos clave: diferencias centrales entre atención y SSM}
La atención accede explícitamente a todo el contexto para cada consulta (query), ofreciendo flexibilidad en el direccionamiento de contenido pero con una matriz de puntuaciones costosa; los SSM comprimen el pasado en un estado fijo, logrando inferencia eficiente aunque puedan perder detalles accesibles aleatoriamente. Una arquitectura híbrida usaría atención en capas locales o críticas, y espacio de estados o convoluciones en la mayoría de las capas, buscando así un punto óptimo (Pareto) específico para la carga de trabajo (workload-specific).
\end{knowledgebox}

Al comparar JEPA, SSM y Transformers también hay que fijar la representación y la tarea. Si JEPA usa un codificador visual más potente mientras que el modelo base predice píxeles directamente, las diferencias de rendimiento no provienen totalmente del objetivo; si un SSM usa un contexto más largo o un tokenizador diferente, las diferencias en throughput no se pueden atribuir directamente a las actualizaciones de estado. Un experimento justo debe alinear parámetros, tokens de entrenamiento, FLOPs, tiempo real (wall-clock), longitud de secuencia y aumento de datos, reportando por separado prellenado (prefill), decodificación, memoria y calidad. Lo más importante es probar los modos de fallo: si la atención es más fuerte en copia exacta y recuperación arbitrada, si el SSM es más estable en secuencias largas en flujo, y si el latente JEPA realmente soporta planificación condicionada a acciones. Solo estos resultados a nivel de carga de trabajo pueden decidir arquitecturas combinadas, no reemplazados por símbolos de complejidad.

\subsection{Resumen de la sección}

JEPA/modelos mundiales cambian el objetivo de aprendizaje; los SSM cambian la ruta de cálculo. El primero busca representaciones más abstractas y predecibles del estado del entorno; el segundo busca eficiencia en secuencias largas. Ninguno es automáticamente superior al Transformer solo por su nombre conceptual; deben compararse bajo condiciones explícitas de tokenización, datos, tarea, cómputo y evaluación.


\section{Síntesis y ampliación}

\subsection{Un hilo conductor sistémico que atraviesa toda la clase}

Esta sesión de repaso puede condensarse en una cadena de dependencias: el mundo original se mapea primero mediante tokenización/encoder a representaciones; la arquitectura determina cómo interactúa la información; los datos de preentrenamiento y el objetivo definen qué patrones estadísticos se forman en los parámetros; el post-entrenamiento y el cómputo en tiempo de inferencia modifican el comportamiento; la recuperación, la memoria, las herramientas y el bucle del agente añaden estado externo; la interpretabilidad, el alineamiento y el protocolo de evidencia deciden si el sistema puede entenderse y restringirse; JEPA y SSM exploran respectivamente nuevos objetos de predicción y rutas de cálculo.

Esta cadena también explica por qué las fallas sistémicas suelen propagarse transversalmente entre capas. El tokenizador al romper los límites de las entidades degrada la consulta de recuperación; si esta devuelve evidencia errónea, la trazabilidad del razonamiento se alarga sobre premisas incorrectas; si el modelo de preferencias favorece un tono confiado, el agente podría reducir la abstención necesaria; si la memoria guarda resúmenes no verificados, cada tarea posterior convertirá un error en un estado persistente. Por el contrario, las reparaciones locales también pueden fracasar: ampliar solo el contexto no corrige los permisos erróneos, aplicar solo DPO no actualiza conocimientos obsoletos y añadir solo un verificador no recupera variables inobservables. La auditoría de ingeniería debe ubicar problemas capa por capa a lo largo de representación, datos, objetivo, estado, acción y evaluación, en lugar de atribuir todos los problemas a que "el modelo no es lo bastante inteligente".

Por ello, un informe de sistema confiable debe proporcionar simultáneamente los límites temporales de los datos de entrenamiento, las versiones del modelo y del tokenizador, los métodos de post-entrenamiento, las herramientas y permisos disponibles, el ciclo de vida de la memoria, los conjuntos de evaluación y las métricas de costo. Si faltan estas informaciones, incluso si se reproducen las puntuaciones individuales de un benchmark, no será posible replicar el comportamiento real del sistema; el nombre del modelo no constituye una condición experimental completa.

De igual manera, cualquier afirmación sobre una "alternativa a Transformer" debe especificar bajo qué longitud de secuencia, lote, hardware, presupuesto de entrenamiento y umbrales de calidad se realizó la comparación. Los rankings arquitectónicos fuera del workload carecen de un significado estable; lo verdaderamente transferible son los métodos de medición y las condiciones de contorno.

\begin{importantbox}{Las diez preguntas más transferibles de toda la clase}
\begin{enumerate}
  \item ¿En qué tokens se corta la entrada y qué información se pierde en la entrada?
  \item ¿Qué patrones estadísticos recompensa la función objetivo y cuáles omite?
  \item ¿Cuáles son el origen de los datos, la mezcla, el orden y la frecuencia de actualización?
  \item ¿Dónde está el cuello de botella computacional: parámetros, atención, memoria, recuperación o herramientas?
  \item ¿El post-entrenamiento modifica pesos, cambia la búsqueda de candidatos o añade estado externo?
  \item ¿Cuáles son las observaciones, acciones, permisos, presupuesto y condiciones de parada del agente?
  \item ¿Qué miden las métricas: pérdida, éxito de tareas, preferencia o impacto real?
  \item ¿De dónde provienen los errores: creencias, observación, planificación o proxy de recompensa?
  \item ¿Cómo guarda, actualiza, resuelve conflictos y olvida el sistema?
  \item ¿Bajo condiciones iguales de datos, cómputo y servicio, qué dimensión del Pareto mejora la nueva arquitectura?
\end{enumerate}
\end{importantbox}

\subsection{Tres límites de evidencia}

Primero, los estudios sobre BabyLM, RAG, CGLS, fMRI y alucinación presentados en clase corresponden a diseños experimentales específicos y no pueden extrapolarse como leyes universales desde una única escala o conjunto de datos. Segundo, las demostraciones de aplicación y los benchmarks no sustituyen el costo, los permisos, el desplazamiento de distribución y la validación externa presentes en un flujo de trabajo real. Tercero, conceptos como world model, aprendizaje continuo y alineamiento deben definir primero qué objeto se actualiza y qué protocolo de evaluación aplica; de lo contrario, la misma palabra cubrirá problemas mutuamente distintos.

\subsection{Preguntas abiertas}

Lo siguiente merece seguir cuestionando: ¿es posible obtener cálculo de secuencias largas lineales sin sacrificar la capacidad de acceso aleatorio? ¿Podemos unificar recuperación, actualización de parámetros y conflicto de memoria en un sistema de aprendizaje continuo verificable? ¿Podemos construir un benchmark del agente con diagnóstico estratificado de las creencias del modelo del mundo, el planificador y la acción? ¿Podemos hacer que la supervisión del proceso y la evidencia mecanicista se validen mutuamente en lugar de recompensar conjuntamente un estilo superficial de razonamiento?

\subsection{Lecturas adicionales}

\begin{itemize}
  \item Curso oficial: \url{https://web.stanford.edu/class/cs25/}
  \item Grabaciones oficiales: \url{https://www.youtube.com/watch?v=bHSDPgZYie0}
  \item Baby Scale: \url{https://arxiv.org/abs/2603.29522}
  \item Bilingual BabyLM: \url{https://arxiv.org/abs/2603.29552}
  \item Preentrenamiento--escalado RAG: \url{https://arxiv.org/abs/2604.00715}
  \item Escalado de capas guiado por currículo: \url{https://arxiv.org/abs/2506.11389}
  \item Representación fMRI cruzada de red: \url{https://arxiv.org/abs/2603.00786}
  \item Definición unificada de alucinación: \url{https://arxiv.org/abs/2512.21577}
\end{itemize}

\end{document}
